
\documentclass{article}
\usepackage{colortbl}
\usepackage{makecell}
\usepackage{multirow}
\usepackage{supertabular}

\begin{document}

\newcounter{utterance}

\centering \large Interaction Transcript for game `clean\_up', experiment `1\_medium\_7obj\_et', episode 3 with nemotron{-}3{-}ultra.
\vspace{24pt}

{ \footnotesize  \setcounter{utterance}{1}
\setlength{\tabcolsep}{0pt}
\begin{supertabular}{c@{$\;$}|p{.15\linewidth}@{}p{.15\linewidth}p{.15\linewidth}p{.15\linewidth}p{.15\linewidth}p{.15\linewidth}}
    \# & $\;$A & \multicolumn{4}{c}{Game Master} & $\;\:$B\\
    \hline

    \theutterance \stepcounter{utterance}  
    & & \multicolumn{4}{p{0.6\linewidth}}{
        \cellcolor[rgb]{0.9,0.9,0.9}{
            \makecell[{{p{\linewidth}}}]{
                \texttt{\tiny{[P1$\langle$GM]}}
                \texttt{Ma olen teie mängujuht ja te mängite ühismängu järgmise ruudustikuga mängulauana:} \\
\\ 
\texttt{```} \\
\texttt{ 1234567} \\
\texttt{╔═╤═════╗} \\
\texttt{║◌│◌A◌◌◌║ 1} \\
\texttt{╟─┤◌T◌◌◌║ 2} \\
\texttt{║O└─┐◌◌◌║ 3} \\
\texttt{║◌◌◌├─┐◌║ 4} \\
\texttt{║◌◌S│◌│◌║ 5} \\
\texttt{╟─┬─┴─┼─╢ 6} \\
\texttt{║P│◌DM│◌║ 7} \\
\texttt{╚═╧═══╧═╝} \\
\\ 
\texttt{```} \\
\\ 
\texttt{* Ülemisel serval on x{-}koordinaadid, mis suurenevad paremale, ja paremal serval y{-}koordinaadid, mis suurenevad allapoole.} \\
\texttt{* Järgmised objektid on teie ruudustikule juhuslikult paigutatud: 'P', 'O', 'T', 'S', 'D', 'A', 'M'.} \\
\\ 
\\ 
\texttt{Teine mängija näeb mängulaua varianti, kus objektid on paigutatud erinevatesse juhuslikesse asukohtadesse. Te ei näe teise mängija mängulauda ja tema ei näe teie mängulauda.} \\
\\ 
\\ 
\texttt{**Eesmärk:**} \\
\\ 
\texttt{Mõlemad mängijad peavad oma taustal olevaid objekte liigutama nii, et identsed objektid satuksid samadesse koordinaatidesse. Te peate teise mängijaga suhtlema, et kokku leppida ühises eesmärgi konfiguratsioonis.} \\
\\ 
\\ 
\texttt{**Reeglid:**} \\
\\ 
\texttt{* Igal käigul saate saata täpselt ühe kahest järgmisest käsust:} \\
\texttt{ 1. `ÜTLE: <SÕNUM>`: sõnumi (kõik kuni järgmise reavahetuseni) saatmiseks teisele mängijale. Ma edastan selle teie partnerile.} \\
\texttt{2. `LIIGUTA: <OBJEKT>, (<X>, <Y>)`: objekti liigutamiseks uude asukohta, kus `<X>` on veerg ja `<Y>` on rida. Ma teatan teile, kas teie käik õnnestus või mitte.} \\
\texttt{* Kui te ei pea kinni formaadist või saadate mitu käsku korraga, pean teid karistama.} \\
\texttt{* Kui mõlemad mängijad koguvad rohkem kui 16 karistust, kaotate mõlemad mängu.} \\
\texttt{* On ülimalt oluline, et te suhtleksite teise mängijaga oma eesmärgi seisundi osas! *Ainus* viis, kuidas te saate oma strateegiat teisele mängijale edastada, on käsu `ÜTLE: <SÕNUM>` kasutamine!} \\
\\ 
\\ 
\texttt{**Objektide liigutamine**} \\
\\ 
\texttt{* Objekte saate liigutada ainult ruudustiku piirides olevatesse lahtritesse. Sihtlahter peab olema tühi, st see peab sisaldama ainult sümbolit '◌'.} \\
\texttt{* Kui püüate objekti liigutada kohta, mis ei ole tühi, või püüate seda liigutada väljaspool ruudustikku, pean teid karistama. Saate uue katse.} \\
\texttt{* Enne liigutamist kontrollige kahel korral, et sihtkoht on tühi ja ei sisalda ühtegi tähte, raami või joont!} \\
\\ 
\\ 
\texttt{**Mängu lõpp**} \\
\\ 
\texttt{Kui arvate, et olete saavutanud kõigi objektide joondamise eesmärgi, võite paluda teisel mängijal mäng lõpetada, saates `ÜTLE: valmis?`. Kui teine mängija palub teil mängu lõpetada ja te vastate `ÜTLE: valmis!`, siis mäng lõpeb.} \\
\\ 
\texttt{Mõlemad mängijad võidavad, kui mäng lõpeb 28 vooru jooksul, kus üks voor on määratletud kahe mängija poolt ühe kehtiva käsu saatmisena.} \\
\\ 
\texttt{**Punktiarvestus:**} \\
\\ 
\texttt{Mida lähemal on identsed objektid mõlemal mängulaual, seda rohkem punkte saate. Karistused vähendavad teie punkte. Kas suudate rekordit ületada?} \\
\\ 
\\ 
\texttt{Palun saada teisele mängijale sõnum, et mängu alustada!} \\
            }
        }
    }
    & \\ \\

    \theutterance \stepcounter{utterance}  
    & \multicolumn{4}{p{0.6\linewidth}}{
        \cellcolor[rgb]{0.9,0.9,0.9}{
            \makecell[{{p{\linewidth}}}]{
                \texttt{\tiny{[P1$\rangle$GM]}}
                \texttt{ÜTLE: Tere! Alustame koos objekti joondamist. Esmalt peaksime kokku lepima, millisesse asendisse iga objekt liigutame. Kuna minu laual on mõned objektid juba üsna hästi paiguti (nt D ja M on reas 7 veergudes 5 ja 6), võiksime proovida joondada kõik read 7. Kas see sobib teile? Palun öelge oma objektide asukohad, et leida ühine laius.} \\
            }
        }
    }
    & & \\ \\

    \theutterance \stepcounter{utterance}  
    & & & \multicolumn{2}{p{0.3\linewidth}}{
        \cellcolor[rgb]{0.9,0.9,0.9}{
            \makecell[{{p{\linewidth}}}]{
                \texttt{\tiny{[GM$|$GM]}}
                \texttt{ÜTLE: Tere! Alustame koos objekti joondamist. Esmalt peaksime kokku lepima, millisesse asendisse iga objekt liigutame. Kuna minu laual on mõned objektid juba üsna hästi paiguti (nt D ja M on reas 7 veergudes 5 ja 6), võiksime proovida joondada kõik read 7. Kas see sobib teile? Palun öelge oma objektide asukohad, et leida ühine laius.} \\
            }
        }
    }
    & & \\ \\

    \theutterance \stepcounter{utterance}  
    & & \multicolumn{4}{p{0.6\linewidth}}{
        \cellcolor[rgb]{0.9,0.9,0.9}{
            \makecell[{{p{\linewidth}}}]{
                \texttt{\tiny{[GM$\rangle$P2]}}
                \texttt{Ma olen teie mängujuht ja te mängite ühismängu järgmise ruudustikuga mängulauana:} \\
\\ 
\texttt{```} \\
\texttt{ 1234567} \\
\texttt{╔═╤═════╗} \\
\texttt{║P│S◌D◌◌║ 1} \\
\texttt{╟─┤◌◌◌◌◌║ 2} \\
\texttt{║◌└─┐◌A◌║ 3} \\
\texttt{║◌◌M├─┐◌║ 4} \\
\texttt{║◌T◌│◌│O║ 5} \\
\texttt{╟─┬─┴─┼─╢ 6} \\
\texttt{║◌│◌◌◌│◌║ 7} \\
\texttt{╚═╧═══╧═╝} \\
\\ 
\texttt{```} \\
\\ 
\texttt{* Ülemisel serval on x{-}koordinaadid, mis suurenevad paremale, ja paremal serval y{-}koordinaadid, mis suurenevad allapoole.} \\
\texttt{* Järgmised objektid on teie ruudustikule juhuslikult paigutatud: 'P', 'O', 'T', 'S', 'D', 'A', 'M'.} \\
\\ 
\\ 
\texttt{Teine mängija näeb mängulaua varianti, kus objektid on paigutatud erinevatesse juhuslikesse asukohtadesse. Te ei näe teise mängija mängulauda ja tema ei näe teie mängulauda.} \\
\\ 
\\ 
\texttt{**Eesmärk:**} \\
\\ 
\texttt{Mõlemad mängijad peavad oma taustal olevaid objekte liigutama nii, et identsed objektid satuksid samadesse koordinaatidesse. Te peate teise mängijaga suhtlema, et kokku leppida ühises eesmärgi konfiguratsioonis.} \\
\\ 
\\ 
\texttt{**Reeglid:**} \\
\\ 
\texttt{* Igal käigul saate saata täpselt ühe kahest järgmisest käsust:} \\
\texttt{ 1. `ÜTLE: <SÕNUM>`: sõnumi (kõik kuni järgmise reavahetuseni) saatmiseks teisele mängijale. Ma edastan selle teie partnerile.} \\
\texttt{2. `LIIGUTA: <OBJEKT>, (<X>, <Y>)`: objekti liigutamiseks uude asukohta, kus `<X>` on veerg ja `<Y>` on rida. Ma teatan teile, kas teie käik õnnestus või mitte.} \\
\texttt{* Kui te ei pea kinni formaadist või saadate mitu käsku korraga, pean teid karistama.} \\
\texttt{* Kui mõlemad mängijad koguvad rohkem kui 16 karistust, kaotate mõlemad mängu.} \\
\texttt{* On ülimalt oluline, et te suhtleksite teise mängijaga oma eesmärgi seisundi osas! *Ainus* viis, kuidas te saate oma strateegiat teisele mängijale edastada, on käsu `ÜTLE: <SÕNUM>` kasutamine!} \\
\\ 
\\ 
\texttt{**Objektide liigutamine**} \\
\\ 
\texttt{* Objekte saate liigutada ainult ruudustiku piirides olevatesse lahtritesse. Sihtlahter peab olema tühi, st see peab sisaldama ainult sümbolit '◌'.} \\
\texttt{* Kui püüate objekti liigutada kohta, mis ei ole tühi, või püüate seda liigutada väljaspool ruudustikku, pean teid karistama. Saate uue katse.} \\
\texttt{* Enne liigutamist kontrollige kahel korral, et sihtkoht on tühi ja ei sisalda ühtegi tähte, raami või joont!} \\
\\ 
\\ 
\texttt{**Mängu lõpp**} \\
\\ 
\texttt{Kui arvate, et olete saavutanud kõigi objektide joondamise eesmärgi, võite paluda teisel mängijal mäng lõpetada, saates `ÜTLE: valmis?`. Kui teine mängija palub teil mängu lõpetada ja te vastate `ÜTLE: valmis!`, siis mäng lõpeb.} \\
\\ 
\texttt{Mõlemad mängijad võidavad, kui mäng lõpeb 28 vooru jooksul, kus üks voor on määratletud kahe mängija poolt ühe kehtiva käsu saatmisena.} \\
\\ 
\texttt{**Punktiarvestus:**} \\
\\ 
\texttt{Mida lähemal on identsed objektid mõlemal mängulaual, seda rohkem punkte saate. Karistused vähendavad teie punkte. Kas suudate rekordit ületada?} \\
\\ 
\texttt{Teine mängija alustas mängu selle sõnumi saatmisega:} \\
\\ 
\texttt{"Tere! Alustame koos objekti joondamist. Esmalt peaksime kokku lepima, millisesse asendisse iga objekt liigutame. Kuna minu laual on mõned objektid juba üsna hästi paiguti (nt D ja M on reas 7 veergudes 5 ja 6), võiksime proovida joondada kõik read 7. Kas see sobib teile? Palun öelge oma objektide asukohad, et leida ühine laius."} \\
\\ 
\texttt{Mis on sinu esimene käsk?} \\
            }
        }
    }
    & \\ \\

    \theutterance \stepcounter{utterance}  
    & & & \multicolumn{4}{p{0.6\linewidth}}{
        \cellcolor[rgb]{0.9,0.9,0.9}{
            \makecell[{{p{\linewidth}}}]{
                \texttt{\tiny{[GM$\langle$P2]}}
                \texttt{ÜTLE: Tere! Minu laual on objekid järgmistes koordinaatides: P(1,1), S(2,1), D(4,1), A(6,3), M(3,4), T(2,5), O(7,5). Rea 7 laius (veergud 1{-}7) tundub hea eesmärk, kuna see mahutab kõik 7 objekti. Eelistan joondada need veergudesse 1–7 vastavalt sorteeritud tähestikulisele järjekorrale (A, D, M, O, P, S, T), et meil oleks deterministicne ja korduva väärtuseta plaan. Seega sihtkoordinaadid: A(1,7), D(2,7), M(3,7), O(4,7), P(5,7), S(6,7), T(7,7). Kas see plaan sobib teile?} \\
            }
        }
    }
     \\ \\

    \theutterance \stepcounter{utterance}  
    & & & \multicolumn{2}{p{0.3\linewidth}}{
        \cellcolor[rgb]{0.9,0.9,0.9}{
            \makecell[{{p{\linewidth}}}]{
                \texttt{\tiny{[GM$|$GM]}}
                \texttt{ÜTLE: Tere! Minu laual on objekid järgmistes koordinaatides: P(1,1), S(2,1), D(4,1), A(6,3), M(3,4), T(2,5), O(7,5). Rea 7 laius (veergud 1{-}7) tundub hea eesmärk, kuna see mahutab kõik 7 objekti. Eelistan joondada need veergudesse 1–7 vastavalt sorteeritud tähestikulisele järjekorrale (A, D, M, O, P, S, T), et meil oleks deterministicne ja korduva väärtuseta plaan. Seega sihtkoordinaadid: A(1,7), D(2,7), M(3,7), O(4,7), P(5,7), S(6,7), T(7,7). Kas see plaan sobib teile?} \\
            }
        }
    }
    & & \\ \\

    \theutterance \stepcounter{utterance}  
    & & \multicolumn{4}{p{0.6\linewidth}}{
        \cellcolor[rgb]{0.9,0.9,0.9}{
            \makecell[{{p{\linewidth}}}]{
                \texttt{\tiny{[P1$\langle$GM]}}
                \texttt{Sinu sõnum on teisele mängijale edastatud.} \\
\\ 
\\ 
\texttt{Sa mängid praegu vooru 1 maksimaalsest 28 voorust.} \\
\texttt{Teil on kollektiivselt kogunenud 0 karistust 16 karistusest.} \\
\texttt{Teine mängija saatis selle sõnumi:} \\
\\ 
\\ 
\texttt{"Tere! Minu laual on objekid järgmistes koordinaatides: P(1,1), S(2,1), D(4,1), A(6,3), M(3,4), T(2,5), O(7,5). Rea 7 laius (veergud 1{-}7) tundub hea eesmärk, kuna see mahutab kõik 7 objekti. Eelistan joondada need veergudesse 1–7 vastavalt sorteeritud tähestikulisele järjekorrale (A, D, M, O, P, S, T), et meil oleks deterministicne ja korduva väärtuseta plaan. Seega sihtkoordinaadid: A(1,7), D(2,7), M(3,7), O(4,7), P(5,7), S(6,7), T(7,7). Kas see plaan sobib teile?"} \\
\\ 
\\ 
\texttt{Mis on sinu järgmine käsk?} \\
            }
        }
    }
    & \\ \\

    \theutterance \stepcounter{utterance}  
    & \multicolumn{4}{p{0.6\linewidth}}{
        \cellcolor[rgb]{0.9,0.9,0.9}{
            \makecell[{{p{\linewidth}}}]{
                \texttt{\tiny{[P1$\rangle$GM]}}
                \texttt{ÜTLE: Plaan on perfektne! Sorteeritud tähestikuline järjekord veergudesse 1{-}7 real 7 tagab täielikult deterministicse ja korduva väärtuseta eesmärgi. Minu praegused asukohad: A(3,1), D(5,7), M(6,7), O(1,3), P(1,7), S(3,5), T(2,2). Alustan kohe liigutamisega. Liigutan esimese objekti, mis on juba õiges veerus: D(5,7) {-}> (2,7).} \\
            }
        }
    }
    & & \\ \\

    \theutterance \stepcounter{utterance}  
    & & & \multicolumn{2}{p{0.3\linewidth}}{
        \cellcolor[rgb]{0.9,0.9,0.9}{
            \makecell[{{p{\linewidth}}}]{
                \texttt{\tiny{[GM$|$GM]}}
                \texttt{ÜTLE: Plaan on perfektne! Sorteeritud tähestikuline järjekord veergudesse 1{-}7 real 7 tagab täielikult deterministicse ja korduva väärtuseta eesmärgi. Minu praegused asukohad: A(3,1), D(5,7), M(6,7), O(1,3), P(1,7), S(3,5), T(2,2). Alustan kohe liigutamisega. Liigutan esimese objekti, mis on juba õiges veerus: D(5,7) {-}> (2,7).} \\
            }
        }
    }
    & & \\ \\

    \theutterance \stepcounter{utterance}  
    & & \multicolumn{4}{p{0.6\linewidth}}{
        \cellcolor[rgb]{0.9,0.9,0.9}{
            \makecell[{{p{\linewidth}}}]{
                \texttt{\tiny{[GM$\rangle$P2]}}
                \texttt{Sinu sõnum on teisele mängijale edastatud.} \\
\\ 
\\ 
\texttt{Sa mängid praegu vooru 2 maksimaalsest 28 voorust.} \\
\texttt{Teil on kollektiivselt kogunenud 0 karistust 16 karistusest.} \\
\texttt{Teine mängija saatis selle sõnumi:} \\
\\ 
\\ 
\texttt{"Plaan on perfektne! Sorteeritud tähestikuline järjekord veergudesse 1{-}7 real 7 tagab täielikult deterministicse ja korduva väärtuseta eesmärgi. Minu praegused asukohad: A(3,1), D(5,7), M(6,7), O(1,3), P(1,7), S(3,5), T(2,2). Alustan kohe liigutamisega. Liigutan esimese objekti, mis on juba õiges veerus: D(5,7) {-}> (2,7)."} \\
\\ 
\\ 
\texttt{Mis on sinu järgmine käsk?} \\
            }
        }
    }
    & \\ \\

    \theutterance \stepcounter{utterance}  
    & & & \multicolumn{4}{p{0.6\linewidth}}{
        \cellcolor[rgb]{0.9,0.9,0.9}{
            \makecell[{{p{\linewidth}}}]{
                \texttt{\tiny{[GM$\langle$P2]}}
                \texttt{LIIGUTA: P, (5, 7)} \\
            }
        }
    }
     \\ \\

    \theutterance \stepcounter{utterance}  
    & & & \multicolumn{2}{p{0.3\linewidth}}{
        \cellcolor[rgb]{0.9,0.9,0.9}{
            \makecell[{{p{\linewidth}}}]{
                \texttt{\tiny{[GM$|$GM]}}
                \texttt{LIIGUTA: P, (5, 7)} \\
            }
        }
    }
    & & \\ \\

    \theutterance \stepcounter{utterance}  
    & & & \multicolumn{2}{p{0.3\linewidth}}{
        \cellcolor[rgb]{0.9,0.9,0.9}{
            \makecell[{{p{\linewidth}}}]{
                \texttt{\tiny{[GM$|$GM]}}
                \texttt{P moved to (5, 7)} \\
            }
        }
    }
    & & \\ \\

    \theutterance \stepcounter{utterance}  
    & & \multicolumn{4}{p{0.6\linewidth}}{
        \cellcolor[rgb]{0.9,0.9,0.9}{
            \makecell[{{p{\linewidth}}}]{
                \texttt{\tiny{[P1$\langle$GM]}}
                \texttt{Sinu sõnum on teisele mängijale edastatud.} \\
\\ 
\\ 
\texttt{Sa mängid praegu vooru 2 maksimaalsest 28 voorust.} \\
\texttt{Teil on kollektiivselt kogunenud 0 karistust 16 karistusest.} \\
\texttt{Teine mängija liigutas objekti oma ruudustikul.} \\
\\ 
\\ 
\texttt{Mis on sinu järgmine käsk?} \\
            }
        }
    }
    & \\ \\

    \theutterance \stepcounter{utterance}  
    & \multicolumn{4}{p{0.6\linewidth}}{
        \cellcolor[rgb]{0.9,0.9,0.9}{
            \makecell[{{p{\linewidth}}}]{
                \texttt{\tiny{[P1$\rangle$GM]}}
                \texttt{LIIGUTA: D, (2, 7)} \\
            }
        }
    }
    & & \\ \\

    \theutterance \stepcounter{utterance}  
    & & & \multicolumn{2}{p{0.3\linewidth}}{
        \cellcolor[rgb]{0.9,0.9,0.9}{
            \makecell[{{p{\linewidth}}}]{
                \texttt{\tiny{[GM$|$GM]}}
                \texttt{LIIGUTA: D, (2, 7)} \\
            }
        }
    }
    & & \\ \\

    \theutterance \stepcounter{utterance}  
    & & & \multicolumn{2}{p{0.3\linewidth}}{
        \cellcolor[rgb]{0.9,0.9,0.9}{
            \makecell[{{p{\linewidth}}}]{
                \texttt{\tiny{[GM$|$GM]}}
                \texttt{Karistus: (2,7) ei ole tühi, vaid sisaldab objekti '│'.} \\
\\ 
\texttt{Teil on kollektiivselt kogunenud 1 karistust 16 karistusest.} \\
\texttt{Palun proovi uuesti!} \\
            }
        }
    }
    & & \\ \\

    \theutterance \stepcounter{utterance}  
    & & \multicolumn{4}{p{0.6\linewidth}}{
        \cellcolor[rgb]{0.9,0.9,0.9}{
            \makecell[{{p{\linewidth}}}]{
                \texttt{\tiny{[P1$\langle$GM]}}
                \texttt{Karistus: (2,7) ei ole tühi, vaid sisaldab objekti '│'.} \\
\\ 
\texttt{Teil on kollektiivselt kogunenud 1 karistust 16 karistusest.} \\
\texttt{Palun proovi uuesti!} \\
            }
        }
    }
    & \\ \\

    \theutterance \stepcounter{utterance}  
    & \multicolumn{4}{p{0.6\linewidth}}{
        \cellcolor[rgb]{0.9,0.9,0.9}{
            \makecell[{{p{\linewidth}}}]{
                \texttt{\tiny{[P1$\rangle$GM]}}
                \texttt{ÜTLE: Saan karistuse {-} (2,7) on seina osa. Minu ruudustikus on rea 7 veergud 2{-}4 seinaelementid. Kas teie real 7 on veergud 1{-}7 kõik vabad (ainult '◌')? Kui jah, liigutan D sinna (1,7) ja teised vastavalt. Kui ei, peaksime valima teise rea (nt rida 4 või 5). Palun teavitage.} \\
            }
        }
    }
    & & \\ \\

    \theutterance \stepcounter{utterance}  
    & & & \multicolumn{2}{p{0.3\linewidth}}{
        \cellcolor[rgb]{0.9,0.9,0.9}{
            \makecell[{{p{\linewidth}}}]{
                \texttt{\tiny{[GM$|$GM]}}
                \texttt{ÜTLE: Saan karistuse {-} (2,7) on seina osa. Minu ruudustikus on rea 7 veergud 2{-}4 seinaelementid. Kas teie real 7 on veergud 1{-}7 kõik vabad (ainult '◌')? Kui jah, liigutan D sinna (1,7) ja teised vastavalt. Kui ei, peaksime valima teise rea (nt rida 4 või 5). Palun teavitage.} \\
            }
        }
    }
    & & \\ \\

    \theutterance \stepcounter{utterance}  
    & & \multicolumn{4}{p{0.6\linewidth}}{
        \cellcolor[rgb]{0.9,0.9,0.9}{
            \makecell[{{p{\linewidth}}}]{
                \texttt{\tiny{[GM$\rangle$P2]}}
                \texttt{Edukalt teisaldatud 'P' asukohta (5,7). Teie uuendatud ruudustik näeb välja selline:} \\
\texttt{```} \\
\texttt{ 1234567} \\
\texttt{╔═╤═════╗} \\
\texttt{║◌│S◌D◌◌║ 1} \\
\texttt{╟─┤◌◌◌◌◌║ 2} \\
\texttt{║◌└─┐◌A◌║ 3} \\
\texttt{║◌◌M├─┐◌║ 4} \\
\texttt{║◌T◌│◌│O║ 5} \\
\texttt{╟─┬─┴─┼─╢ 6} \\
\texttt{║◌│◌◌P│◌║ 7} \\
\texttt{╚═╧═══╧═╝} \\
\texttt{```} \\
\texttt{Sa mängid praegu vooru 3 maksimaalsest 28 voorust.} \\
\texttt{Teil on kollektiivselt kogunenud 1 karistust 16 karistusest.} \\
\texttt{Teine mängija saatis selle sõnumi:} \\
\\ 
\\ 
\texttt{"Saan karistuse {-} (2,7) on seina osa. Minu ruudustikus on rea 7 veergud 2{-}4 seinaelementid. Kas teie real 7 on veergud 1{-}7 kõik vabad (ainult '◌')? Kui jah, liigutan D sinna (1,7) ja teised vastavalt. Kui ei, peaksime valima teise rea (nt rida 4 või 5). Palun teavitage."} \\
\\ 
\\ 
\texttt{Mis on sinu järgmine käsk?} \\
            }
        }
    }
    & \\ \\

    \theutterance \stepcounter{utterance}  
    & & & \multicolumn{4}{p{0.6\linewidth}}{
        \cellcolor[rgb]{0.9,0.9,0.9}{
            \makecell[{{p{\linewidth}}}]{
                \texttt{\tiny{[GM$\langle$P2]}}
                \texttt{ÜTLE: Minu rida 7 on täiesti vaba (veergud 1{-}7 on kõik '◌', seinaid pole). Jätkame plaani: A(1,7), D(2,7), M(3,7), O(4,7), P(5,7), S(6,7), T(7,7). Teie järgmine käik: liiguta D (1,7). Minu järgmine käik: liigutan A (1,7).} \\
            }
        }
    }
     \\ \\

    \theutterance \stepcounter{utterance}  
    & & & \multicolumn{2}{p{0.3\linewidth}}{
        \cellcolor[rgb]{0.9,0.9,0.9}{
            \makecell[{{p{\linewidth}}}]{
                \texttt{\tiny{[GM$|$GM]}}
                \texttt{ÜTLE: Minu rida 7 on täiesti vaba (veergud 1{-}7 on kõik '◌', seinaid pole). Jätkame plaani: A(1,7), D(2,7), M(3,7), O(4,7), P(5,7), S(6,7), T(7,7). Teie järgmine käik: liiguta D (1,7). Minu järgmine käik: liigutan A (1,7).} \\
            }
        }
    }
    & & \\ \\

    \theutterance \stepcounter{utterance}  
    & & \multicolumn{4}{p{0.6\linewidth}}{
        \cellcolor[rgb]{0.9,0.9,0.9}{
            \makecell[{{p{\linewidth}}}]{
                \texttt{\tiny{[P1$\langle$GM]}}
                \texttt{Sinu sõnum on teisele mängijale edastatud.} \\
\\ 
\\ 
\texttt{Sa mängid praegu vooru 3 maksimaalsest 28 voorust.} \\
\texttt{Teil on kollektiivselt kogunenud 1 karistust 16 karistusest.} \\
\texttt{Teine mängija saatis selle sõnumi:} \\
\\ 
\\ 
\texttt{"Minu rida 7 on täiesti vaba (veergud 1{-}7 on kõik '◌', seinaid pole). Jätkame plaani: A(1,7), D(2,7), M(3,7), O(4,7), P(5,7), S(6,7), T(7,7). Teie järgmine käik: liiguta D (1,7). Minu järgmine käik: liigutan A (1,7)."} \\
\\ 
\\ 
\texttt{Mis on sinu järgmine käsk?} \\
            }
        }
    }
    & \\ \\

    \theutterance \stepcounter{utterance}  
    & \multicolumn{4}{p{0.6\linewidth}}{
        \cellcolor[rgb]{0.9,0.9,0.9}{
            \makecell[{{p{\linewidth}}}]{
                \texttt{\tiny{[P1$\rangle$GM]}}
                \texttt{LIIGUTA: D, (1, 7)} \\
            }
        }
    }
    & & \\ \\

    \theutterance \stepcounter{utterance}  
    & & & \multicolumn{2}{p{0.3\linewidth}}{
        \cellcolor[rgb]{0.9,0.9,0.9}{
            \makecell[{{p{\linewidth}}}]{
                \texttt{\tiny{[GM$|$GM]}}
                \texttt{LIIGUTA: D, (1, 7)} \\
            }
        }
    }
    & & \\ \\

    \theutterance \stepcounter{utterance}  
    & & & \multicolumn{2}{p{0.3\linewidth}}{
        \cellcolor[rgb]{0.9,0.9,0.9}{
            \makecell[{{p{\linewidth}}}]{
                \texttt{\tiny{[GM$|$GM]}}
                \texttt{Karistus: (1,7) ei ole tühi, vaid sisaldab objekti 'P'.} \\
\\ 
\texttt{Teil on kollektiivselt kogunenud 2 karistust 16 karistusest.} \\
\texttt{Palun proovi uuesti!} \\
            }
        }
    }
    & & \\ \\

    \theutterance \stepcounter{utterance}  
    & & \multicolumn{4}{p{0.6\linewidth}}{
        \cellcolor[rgb]{0.9,0.9,0.9}{
            \makecell[{{p{\linewidth}}}]{
                \texttt{\tiny{[P1$\langle$GM]}}
                \texttt{Karistus: (1,7) ei ole tühi, vaid sisaldab objekti 'P'.} \\
\\ 
\texttt{Teil on kollektiivselt kogunenud 2 karistust 16 karistusest.} \\
\texttt{Palun proovi uuesti!} \\
            }
        }
    }
    & \\ \\

    \theutterance \stepcounter{utterance}  
    & \multicolumn{4}{p{0.6\linewidth}}{
        \cellcolor[rgb]{0.9,0.9,0.9}{
            \makecell[{{p{\linewidth}}}]{
                \texttt{\tiny{[P1$\rangle$GM]}}
                \texttt{ÜTLE: Vabandust, unustasin, et minu P on juba (1,7). Liigutan P esmalt ära. Sihtkoht P on (5,7), mis peaks olema vaba. Liigutan P sinna.} \\
            }
        }
    }
    & & \\ \\

    \theutterance \stepcounter{utterance}  
    & & & \multicolumn{2}{p{0.3\linewidth}}{
        \cellcolor[rgb]{0.9,0.9,0.9}{
            \makecell[{{p{\linewidth}}}]{
                \texttt{\tiny{[GM$|$GM]}}
                \texttt{ÜTLE: Vabandust, unustasin, et minu P on juba (1,7). Liigutan P esmalt ära. Sihtkoht P on (5,7), mis peaks olema vaba. Liigutan P sinna.} \\
            }
        }
    }
    & & \\ \\

    \theutterance \stepcounter{utterance}  
    & & \multicolumn{4}{p{0.6\linewidth}}{
        \cellcolor[rgb]{0.9,0.9,0.9}{
            \makecell[{{p{\linewidth}}}]{
                \texttt{\tiny{[GM$\rangle$P2]}}
                \texttt{Sinu sõnum on teisele mängijale edastatud.} \\
\\ 
\\ 
\texttt{Sa mängid praegu vooru 4 maksimaalsest 28 voorust.} \\
\texttt{Teil on kollektiivselt kogunenud 2 karistust 16 karistusest.} \\
\texttt{Teine mängija saatis selle sõnumi:} \\
\\ 
\\ 
\texttt{"Vabandust, unustasin, et minu P on juba (1,7). Liigutan P esmalt ära. Sihtkoht P on (5,7), mis peaks olema vaba. Liigutan P sinna."} \\
\\ 
\\ 
\texttt{Mis on sinu järgmine käsk?} \\
            }
        }
    }
    & \\ \\

    \theutterance \stepcounter{utterance}  
    & & & \multicolumn{4}{p{0.6\linewidth}}{
        \cellcolor[rgb]{0.9,0.9,0.9}{
            \makecell[{{p{\linewidth}}}]{
                \texttt{\tiny{[GM$\langle$P2]}}
                \texttt{LIIGUTA: A, (1, 7)} \\
            }
        }
    }
     \\ \\

    \theutterance \stepcounter{utterance}  
    & & & \multicolumn{2}{p{0.3\linewidth}}{
        \cellcolor[rgb]{0.9,0.9,0.9}{
            \makecell[{{p{\linewidth}}}]{
                \texttt{\tiny{[GM$|$GM]}}
                \texttt{LIIGUTA: A, (1, 7)} \\
            }
        }
    }
    & & \\ \\

    \theutterance \stepcounter{utterance}  
    & & & \multicolumn{2}{p{0.3\linewidth}}{
        \cellcolor[rgb]{0.9,0.9,0.9}{
            \makecell[{{p{\linewidth}}}]{
                \texttt{\tiny{[GM$|$GM]}}
                \texttt{A moved to (1, 7)} \\
            }
        }
    }
    & & \\ \\

    \theutterance \stepcounter{utterance}  
    & & \multicolumn{4}{p{0.6\linewidth}}{
        \cellcolor[rgb]{0.9,0.9,0.9}{
            \makecell[{{p{\linewidth}}}]{
                \texttt{\tiny{[P1$\langle$GM]}}
                \texttt{Sinu sõnum on teisele mängijale edastatud.} \\
\\ 
\\ 
\texttt{Sa mängid praegu vooru 4 maksimaalsest 28 voorust.} \\
\texttt{Teil on kollektiivselt kogunenud 2 karistust 16 karistusest.} \\
\texttt{Teine mängija liigutas objekti oma ruudustikul.} \\
\\ 
\\ 
\texttt{Mis on sinu järgmine käsk?} \\
            }
        }
    }
    & \\ \\

    \theutterance \stepcounter{utterance}  
    & \multicolumn{4}{p{0.6\linewidth}}{
        \cellcolor[rgb]{0.9,0.9,0.9}{
            \makecell[{{p{\linewidth}}}]{
                \texttt{\tiny{[P1$\rangle$GM]}}
                \texttt{LIIGUTA: P, (5, 7)} \\
            }
        }
    }
    & & \\ \\

    \theutterance \stepcounter{utterance}  
    & & & \multicolumn{2}{p{0.3\linewidth}}{
        \cellcolor[rgb]{0.9,0.9,0.9}{
            \makecell[{{p{\linewidth}}}]{
                \texttt{\tiny{[GM$|$GM]}}
                \texttt{LIIGUTA: P, (5, 7)} \\
            }
        }
    }
    & & \\ \\

    \theutterance \stepcounter{utterance}  
    & & & \multicolumn{2}{p{0.3\linewidth}}{
        \cellcolor[rgb]{0.9,0.9,0.9}{
            \makecell[{{p{\linewidth}}}]{
                \texttt{\tiny{[GM$|$GM]}}
                \texttt{Karistus: (5,7) ei ole tühi, vaid sisaldab objekti 'M'.} \\
\\ 
\texttt{Teil on kollektiivselt kogunenud 3 karistust 16 karistusest.} \\
\texttt{Palun proovi uuesti!} \\
            }
        }
    }
    & & \\ \\

    \theutterance \stepcounter{utterance}  
    & & \multicolumn{4}{p{0.6\linewidth}}{
        \cellcolor[rgb]{0.9,0.9,0.9}{
            \makecell[{{p{\linewidth}}}]{
                \texttt{\tiny{[P1$\langle$GM]}}
                \texttt{Karistus: (5,7) ei ole tühi, vaid sisaldab objekti 'M'.} \\
\\ 
\texttt{Teil on kollektiivselt kogunenud 3 karistust 16 karistusest.} \\
\texttt{Palun proovi uuesti!} \\
            }
        }
    }
    & \\ \\

    \theutterance \stepcounter{utterance}  
    & \multicolumn{4}{p{0.6\linewidth}}{
        \cellcolor[rgb]{0.9,0.9,0.9}{
            \makecell[{{p{\linewidth}}}]{
                \texttt{\tiny{[P1$\rangle$GM]}}
                \texttt{ÜTLE: Minu M on algselt (6,7), seega (5,7) on minu laual hõivatud. Minu rida 7: veerg 1=P, 2=seina, 3=seina, 4=seina, 5=M, 6=D, 7=tühi. Vabad veerud real 7 on ainult 7. Teie laual on 1{-}7 vabad. Peaksime koordineerima: kas te liigutate oma M (3,7) enne minema, et vabastada (5,7) mulle? Või kas ma liigutan minu M (6,7) {-}> (7,7) enne P liigutamist? Palun öelge, milline järjekord sobib teile.} \\
            }
        }
    }
    & & \\ \\

    \theutterance \stepcounter{utterance}  
    & & & \multicolumn{2}{p{0.3\linewidth}}{
        \cellcolor[rgb]{0.9,0.9,0.9}{
            \makecell[{{p{\linewidth}}}]{
                \texttt{\tiny{[GM$|$GM]}}
                \texttt{ÜTLE: Minu M on algselt (6,7), seega (5,7) on minu laual hõivatud. Minu rida 7: veerg 1=P, 2=seina, 3=seina, 4=seina, 5=M, 6=D, 7=tühi. Vabad veerud real 7 on ainult 7. Teie laual on 1{-}7 vabad. Peaksime koordineerima: kas te liigutate oma M (3,7) enne minema, et vabastada (5,7) mulle? Või kas ma liigutan minu M (6,7) {-}> (7,7) enne P liigutamist? Palun öelge, milline järjekord sobib teile.} \\
            }
        }
    }
    & & \\ \\

    \theutterance \stepcounter{utterance}  
    & & \multicolumn{4}{p{0.6\linewidth}}{
        \cellcolor[rgb]{0.9,0.9,0.9}{
            \makecell[{{p{\linewidth}}}]{
                \texttt{\tiny{[GM$\rangle$P2]}}
                \texttt{Edukalt teisaldatud 'A' asukohta (1,7). Teie uuendatud ruudustik näeb välja selline:} \\
\texttt{```} \\
\texttt{ 1234567} \\
\texttt{╔═╤═════╗} \\
\texttt{║◌│S◌D◌◌║ 1} \\
\texttt{╟─┤◌◌◌◌◌║ 2} \\
\texttt{║◌└─┐◌◌◌║ 3} \\
\texttt{║◌◌M├─┐◌║ 4} \\
\texttt{║◌T◌│◌│O║ 5} \\
\texttt{╟─┬─┴─┼─╢ 6} \\
\texttt{║A│◌◌P│◌║ 7} \\
\texttt{╚═╧═══╧═╝} \\
\texttt{```} \\
\texttt{Sa mängid praegu vooru 5 maksimaalsest 28 voorust.} \\
\texttt{Teil on kollektiivselt kogunenud 3 karistust 16 karistusest.} \\
\texttt{Teine mängija saatis selle sõnumi:} \\
\\ 
\\ 
\texttt{"Minu M on algselt (6,7), seega (5,7) on minu laual hõivatud. Minu rida 7: veerg 1=P, 2=seina, 3=seina, 4=seina, 5=M, 6=D, 7=tühi. Vabad veerud real 7 on ainult 7. Teie laual on 1{-}7 vabad. Peaksime koordineerima: kas te liigutate oma M (3,7) enne minema, et vabastada (5,7) mulle? Või kas ma liigutan minu M (6,7) {-}> (7,7) enne P liigutamist? Palun öelge, milline järjekord sobib teile."} \\
\\ 
\\ 
\texttt{Mis on sinu järgmine käsk?} \\
            }
        }
    }
    & \\ \\

    \theutterance \stepcounter{utterance}  
    & & & \multicolumn{4}{p{0.6\linewidth}}{
        \cellcolor[rgb]{0.9,0.9,0.9}{
            \makecell[{{p{\linewidth}}}]{
                \texttt{\tiny{[GM$\langle$P2]}}
                \texttt{ÜTLE: Minu M on algselt (3,4), seega mu real 7 on kõik veerud 1{-}7 vabad. Teie plaan: liiguta M (6,7) {-}> (7,7), seejärel P (1,7) {-}> (5,7). Minu järgmiseks liigutan D (4,1) {-}> (2,7).} \\
            }
        }
    }
     \\ \\

    \theutterance \stepcounter{utterance}  
    & & & \multicolumn{2}{p{0.3\linewidth}}{
        \cellcolor[rgb]{0.9,0.9,0.9}{
            \makecell[{{p{\linewidth}}}]{
                \texttt{\tiny{[GM$|$GM]}}
                \texttt{ÜTLE: Minu M on algselt (3,4), seega mu real 7 on kõik veerud 1{-}7 vabad. Teie plaan: liiguta M (6,7) {-}> (7,7), seejärel P (1,7) {-}> (5,7). Minu järgmiseks liigutan D (4,1) {-}> (2,7).} \\
            }
        }
    }
    & & \\ \\

    \theutterance \stepcounter{utterance}  
    & & \multicolumn{4}{p{0.6\linewidth}}{
        \cellcolor[rgb]{0.9,0.9,0.9}{
            \makecell[{{p{\linewidth}}}]{
                \texttt{\tiny{[P1$\langle$GM]}}
                \texttt{Sinu sõnum on teisele mängijale edastatud.} \\
\\ 
\\ 
\texttt{Sa mängid praegu vooru 5 maksimaalsest 28 voorust.} \\
\texttt{Teil on kollektiivselt kogunenud 3 karistust 16 karistusest.} \\
\texttt{Teine mängija saatis selle sõnumi:} \\
\\ 
\\ 
\texttt{"Minu M on algselt (3,4), seega mu real 7 on kõik veerud 1{-}7 vabad. Teie plaan: liiguta M (6,7) {-}> (7,7), seejärel P (1,7) {-}> (5,7). Minu järgmiseks liigutan D (4,1) {-}> (2,7)."} \\
\\ 
\\ 
\texttt{Mis on sinu järgmine käsk?} \\
            }
        }
    }
    & \\ \\

    \theutterance \stepcounter{utterance}  
    & \multicolumn{4}{p{0.6\linewidth}}{
        \cellcolor[rgb]{0.9,0.9,0.9}{
            \makecell[{{p{\linewidth}}}]{
                \texttt{\tiny{[P1$\rangle$GM]}}
                \texttt{LIIGUTA: M, (7, 7)} \\
            }
        }
    }
    & & \\ \\

    \theutterance \stepcounter{utterance}  
    & & & \multicolumn{2}{p{0.3\linewidth}}{
        \cellcolor[rgb]{0.9,0.9,0.9}{
            \makecell[{{p{\linewidth}}}]{
                \texttt{\tiny{[GM$|$GM]}}
                \texttt{LIIGUTA: M, (7, 7)} \\
            }
        }
    }
    & & \\ \\

    \theutterance \stepcounter{utterance}  
    & & & \multicolumn{2}{p{0.3\linewidth}}{
        \cellcolor[rgb]{0.9,0.9,0.9}{
            \makecell[{{p{\linewidth}}}]{
                \texttt{\tiny{[GM$|$GM]}}
                \texttt{M moved to (7, 7)} \\
            }
        }
    }
    & & \\ \\

    \theutterance \stepcounter{utterance}  
    & & \multicolumn{4}{p{0.6\linewidth}}{
        \cellcolor[rgb]{0.9,0.9,0.9}{
            \makecell[{{p{\linewidth}}}]{
                \texttt{\tiny{[GM$\rangle$P2]}}
                \texttt{Sinu sõnum on teisele mängijale edastatud.} \\
\\ 
\\ 
\texttt{Sa mängid praegu vooru 6 maksimaalsest 28 voorust.} \\
\texttt{Teil on kollektiivselt kogunenud 3 karistust 16 karistusest.} \\
\texttt{Teine mängija liigutas objekti oma ruudustikul.} \\
\\ 
\\ 
\texttt{Mis on sinu järgmine käsk?} \\
            }
        }
    }
    & \\ \\

    \theutterance \stepcounter{utterance}  
    & & & \multicolumn{4}{p{0.6\linewidth}}{
        \cellcolor[rgb]{0.9,0.9,0.9}{
            \makecell[{{p{\linewidth}}}]{
                \texttt{\tiny{[GM$\langle$P2]}}
                \texttt{LIIGUTA: D, (2, 7)} \\
            }
        }
    }
     \\ \\

    \theutterance \stepcounter{utterance}  
    & & & \multicolumn{2}{p{0.3\linewidth}}{
        \cellcolor[rgb]{0.9,0.9,0.9}{
            \makecell[{{p{\linewidth}}}]{
                \texttt{\tiny{[GM$|$GM]}}
                \texttt{LIIGUTA: D, (2, 7)} \\
            }
        }
    }
    & & \\ \\

    \theutterance \stepcounter{utterance}  
    & & & \multicolumn{2}{p{0.3\linewidth}}{
        \cellcolor[rgb]{0.9,0.9,0.9}{
            \makecell[{{p{\linewidth}}}]{
                \texttt{\tiny{[GM$|$GM]}}
                \texttt{Karistus: (2,7) ei ole tühi, vaid sisaldab objekti '│'.} \\
\\ 
\texttt{Teil on kollektiivselt kogunenud 4 karistust 16 karistusest.} \\
\texttt{Palun proovi uuesti!} \\
            }
        }
    }
    & & \\ \\

    \theutterance \stepcounter{utterance}  
    & & \multicolumn{4}{p{0.6\linewidth}}{
        \cellcolor[rgb]{0.9,0.9,0.9}{
            \makecell[{{p{\linewidth}}}]{
                \texttt{\tiny{[GM$\rangle$P2]}}
                \texttt{Karistus: (2,7) ei ole tühi, vaid sisaldab objekti '│'.} \\
\\ 
\texttt{Teil on kollektiivselt kogunenud 4 karistust 16 karistusest.} \\
\texttt{Palun proovi uuesti!} \\
            }
        }
    }
    & \\ \\

    \theutterance \stepcounter{utterance}  
    & & & \multicolumn{4}{p{0.6\linewidth}}{
        \cellcolor[rgb]{0.9,0.9,0.9}{
            \makecell[{{p{\linewidth}}}]{
                \texttt{\tiny{[GM$\langle$P2]}}
                \texttt{LIIGUTA: D, (3, 7)} \\
            }
        }
    }
     \\ \\

    \theutterance \stepcounter{utterance}  
    & & & \multicolumn{2}{p{0.3\linewidth}}{
        \cellcolor[rgb]{0.9,0.9,0.9}{
            \makecell[{{p{\linewidth}}}]{
                \texttt{\tiny{[GM$|$GM]}}
                \texttt{LIIGUTA: D, (3, 7)} \\
            }
        }
    }
    & & \\ \\

    \theutterance \stepcounter{utterance}  
    & & & \multicolumn{2}{p{0.3\linewidth}}{
        \cellcolor[rgb]{0.9,0.9,0.9}{
            \makecell[{{p{\linewidth}}}]{
                \texttt{\tiny{[GM$|$GM]}}
                \texttt{D moved to (3, 7)} \\
            }
        }
    }
    & & \\ \\

    \theutterance \stepcounter{utterance}  
    & & \multicolumn{4}{p{0.6\linewidth}}{
        \cellcolor[rgb]{0.9,0.9,0.9}{
            \makecell[{{p{\linewidth}}}]{
                \texttt{\tiny{[P1$\langle$GM]}}
                \texttt{Edukalt teisaldatud 'M' asukohta (7,7). Teie uuendatud ruudustik näeb välja selline:} \\
\texttt{```} \\
\texttt{ 1234567} \\
\texttt{╔═╤═════╗} \\
\texttt{║◌│◌A◌◌◌║ 1} \\
\texttt{╟─┤◌T◌◌◌║ 2} \\
\texttt{║O└─┐◌◌◌║ 3} \\
\texttt{║◌◌◌├─┐◌║ 4} \\
\texttt{║◌◌S│◌│◌║ 5} \\
\texttt{╟─┬─┴─┼─╢ 6} \\
\texttt{║P│◌D◌│M║ 7} \\
\texttt{╚═╧═══╧═╝} \\
\texttt{```} \\
\texttt{Sa mängid praegu vooru 6 maksimaalsest 28 voorust.} \\
\texttt{Teil on kollektiivselt kogunenud 4 karistust 16 karistusest.} \\
\texttt{Teine mängija liigutas objekti oma ruudustikul.} \\
\\ 
\\ 
\texttt{Mis on sinu järgmine käsk?} \\
            }
        }
    }
    & \\ \\

    \theutterance \stepcounter{utterance}  
    & \multicolumn{4}{p{0.6\linewidth}}{
        \cellcolor[rgb]{0.9,0.9,0.9}{
            \makecell[{{p{\linewidth}}}]{
                \texttt{\tiny{[P1$\rangle$GM]}}
                \texttt{LIIGUTA: P, (5, 7)} \\
            }
        }
    }
    & & \\ \\

    \theutterance \stepcounter{utterance}  
    & & & \multicolumn{2}{p{0.3\linewidth}}{
        \cellcolor[rgb]{0.9,0.9,0.9}{
            \makecell[{{p{\linewidth}}}]{
                \texttt{\tiny{[GM$|$GM]}}
                \texttt{LIIGUTA: P, (5, 7)} \\
            }
        }
    }
    & & \\ \\

    \theutterance \stepcounter{utterance}  
    & & & \multicolumn{2}{p{0.3\linewidth}}{
        \cellcolor[rgb]{0.9,0.9,0.9}{
            \makecell[{{p{\linewidth}}}]{
                \texttt{\tiny{[GM$|$GM]}}
                \texttt{P moved to (5, 7)} \\
            }
        }
    }
    & & \\ \\

    \theutterance \stepcounter{utterance}  
    & & \multicolumn{4}{p{0.6\linewidth}}{
        \cellcolor[rgb]{0.9,0.9,0.9}{
            \makecell[{{p{\linewidth}}}]{
                \texttt{\tiny{[GM$\rangle$P2]}}
                \texttt{Edukalt teisaldatud 'D' asukohta (3,7). Teie uuendatud ruudustik näeb välja selline:} \\
\texttt{```} \\
\texttt{ 1234567} \\
\texttt{╔═╤═════╗} \\
\texttt{║◌│S◌◌◌◌║ 1} \\
\texttt{╟─┤◌◌◌◌◌║ 2} \\
\texttt{║◌└─┐◌◌◌║ 3} \\
\texttt{║◌◌M├─┐◌║ 4} \\
\texttt{║◌T◌│◌│O║ 5} \\
\texttt{╟─┬─┴─┼─╢ 6} \\
\texttt{║A│D◌P│◌║ 7} \\
\texttt{╚═╧═══╧═╝} \\
\texttt{```} \\
\texttt{Sa mängid praegu vooru 7 maksimaalsest 28 voorust.} \\
\texttt{Teil on kollektiivselt kogunenud 4 karistust 16 karistusest.} \\
\texttt{Teine mängija liigutas objekti oma ruudustikul.} \\
\\ 
\\ 
\texttt{Mis on sinu järgmine käsk?} \\
            }
        }
    }
    & \\ \\

    \theutterance \stepcounter{utterance}  
    & & & \multicolumn{4}{p{0.6\linewidth}}{
        \cellcolor[rgb]{0.9,0.9,0.9}{
            \makecell[{{p{\linewidth}}}]{
                \texttt{\tiny{[GM$\langle$P2]}}
                \texttt{LIIGUTA: M, (3, 7)} \\
            }
        }
    }
     \\ \\

    \theutterance \stepcounter{utterance}  
    & & & \multicolumn{2}{p{0.3\linewidth}}{
        \cellcolor[rgb]{0.9,0.9,0.9}{
            \makecell[{{p{\linewidth}}}]{
                \texttt{\tiny{[GM$|$GM]}}
                \texttt{LIIGUTA: M, (3, 7)} \\
            }
        }
    }
    & & \\ \\

    \theutterance \stepcounter{utterance}  
    & & & \multicolumn{2}{p{0.3\linewidth}}{
        \cellcolor[rgb]{0.9,0.9,0.9}{
            \makecell[{{p{\linewidth}}}]{
                \texttt{\tiny{[GM$|$GM]}}
                \texttt{Karistus: (3,7) ei ole tühi, vaid sisaldab objekti 'D'.} \\
\\ 
\texttt{Teil on kollektiivselt kogunenud 5 karistust 16 karistusest.} \\
\texttt{Palun proovi uuesti!} \\
            }
        }
    }
    & & \\ \\

    \theutterance \stepcounter{utterance}  
    & & \multicolumn{4}{p{0.6\linewidth}}{
        \cellcolor[rgb]{0.9,0.9,0.9}{
            \makecell[{{p{\linewidth}}}]{
                \texttt{\tiny{[GM$\rangle$P2]}}
                \texttt{Karistus: (3,7) ei ole tühi, vaid sisaldab objekti 'D'.} \\
\\ 
\texttt{Teil on kollektiivselt kogunenud 5 karistust 16 karistusest.} \\
\texttt{Palun proovi uuesti!} \\
            }
        }
    }
    & \\ \\

    \theutterance \stepcounter{utterance}  
    & & & \multicolumn{4}{p{0.6\linewidth}}{
        \cellcolor[rgb]{0.9,0.9,0.9}{
            \makecell[{{p{\linewidth}}}]{
                \texttt{\tiny{[GM$\langle$P2]}}
                \texttt{LIIGUTA: M, (4, 7)} \\
            }
        }
    }
     \\ \\

    \theutterance \stepcounter{utterance}  
    & & & \multicolumn{2}{p{0.3\linewidth}}{
        \cellcolor[rgb]{0.9,0.9,0.9}{
            \makecell[{{p{\linewidth}}}]{
                \texttt{\tiny{[GM$|$GM]}}
                \texttt{LIIGUTA: M, (4, 7)} \\
            }
        }
    }
    & & \\ \\

    \theutterance \stepcounter{utterance}  
    & & & \multicolumn{2}{p{0.3\linewidth}}{
        \cellcolor[rgb]{0.9,0.9,0.9}{
            \makecell[{{p{\linewidth}}}]{
                \texttt{\tiny{[GM$|$GM]}}
                \texttt{M moved to (4, 7)} \\
            }
        }
    }
    & & \\ \\

    \theutterance \stepcounter{utterance}  
    & & \multicolumn{4}{p{0.6\linewidth}}{
        \cellcolor[rgb]{0.9,0.9,0.9}{
            \makecell[{{p{\linewidth}}}]{
                \texttt{\tiny{[P1$\langle$GM]}}
                \texttt{Edukalt teisaldatud 'P' asukohta (5,7). Teie uuendatud ruudustik näeb välja selline:} \\
\texttt{```} \\
\texttt{ 1234567} \\
\texttt{╔═╤═════╗} \\
\texttt{║◌│◌A◌◌◌║ 1} \\
\texttt{╟─┤◌T◌◌◌║ 2} \\
\texttt{║O└─┐◌◌◌║ 3} \\
\texttt{║◌◌◌├─┐◌║ 4} \\
\texttt{║◌◌S│◌│◌║ 5} \\
\texttt{╟─┬─┴─┼─╢ 6} \\
\texttt{║◌│◌DP│M║ 7} \\
\texttt{╚═╧═══╧═╝} \\
\texttt{```} \\
\texttt{Sa mängid praegu vooru 7 maksimaalsest 28 voorust.} \\
\texttt{Teil on kollektiivselt kogunenud 5 karistust 16 karistusest.} \\
\texttt{Teine mängija liigutas objekti oma ruudustikul.} \\
\\ 
\\ 
\texttt{Mis on sinu järgmine käsk?} \\
            }
        }
    }
    & \\ \\

    \theutterance \stepcounter{utterance}  
    & \multicolumn{4}{p{0.6\linewidth}}{
        \cellcolor[rgb]{0.9,0.9,0.9}{
            \makecell[{{p{\linewidth}}}]{
                \texttt{\tiny{[P1$\rangle$GM]}}
                \texttt{LIIGUTA: A, (1, 7)} \\
            }
        }
    }
    & & \\ \\

    \theutterance \stepcounter{utterance}  
    & & & \multicolumn{2}{p{0.3\linewidth}}{
        \cellcolor[rgb]{0.9,0.9,0.9}{
            \makecell[{{p{\linewidth}}}]{
                \texttt{\tiny{[GM$|$GM]}}
                \texttt{LIIGUTA: A, (1, 7)} \\
            }
        }
    }
    & & \\ \\

    \theutterance \stepcounter{utterance}  
    & & & \multicolumn{2}{p{0.3\linewidth}}{
        \cellcolor[rgb]{0.9,0.9,0.9}{
            \makecell[{{p{\linewidth}}}]{
                \texttt{\tiny{[GM$|$GM]}}
                \texttt{A moved to (1, 7)} \\
            }
        }
    }
    & & \\ \\

    \theutterance \stepcounter{utterance}  
    & & \multicolumn{4}{p{0.6\linewidth}}{
        \cellcolor[rgb]{0.9,0.9,0.9}{
            \makecell[{{p{\linewidth}}}]{
                \texttt{\tiny{[GM$\rangle$P2]}}
                \texttt{Edukalt teisaldatud 'M' asukohta (4,7). Teie uuendatud ruudustik näeb välja selline:} \\
\texttt{```} \\
\texttt{ 1234567} \\
\texttt{╔═╤═════╗} \\
\texttt{║◌│S◌◌◌◌║ 1} \\
\texttt{╟─┤◌◌◌◌◌║ 2} \\
\texttt{║◌└─┐◌◌◌║ 3} \\
\texttt{║◌◌◌├─┐◌║ 4} \\
\texttt{║◌T◌│◌│O║ 5} \\
\texttt{╟─┬─┴─┼─╢ 6} \\
\texttt{║A│DMP│◌║ 7} \\
\texttt{╚═╧═══╧═╝} \\
\texttt{```} \\
\texttt{Sa mängid praegu vooru 8 maksimaalsest 28 voorust.} \\
\texttt{Teil on kollektiivselt kogunenud 5 karistust 16 karistusest.} \\
\texttt{Teine mängija liigutas objekti oma ruudustikul.} \\
\\ 
\\ 
\texttt{Mis on sinu järgmine käsk?} \\
            }
        }
    }
    & \\ \\

    \theutterance \stepcounter{utterance}  
    & & & \multicolumn{4}{p{0.6\linewidth}}{
        \cellcolor[rgb]{0.9,0.9,0.9}{
            \makecell[{{p{\linewidth}}}]{
                \texttt{\tiny{[GM$\langle$P2]}}
                \texttt{LIIGUTA: O, (4, 7)} \\
            }
        }
    }
     \\ \\

    \theutterance \stepcounter{utterance}  
    & & & \multicolumn{2}{p{0.3\linewidth}}{
        \cellcolor[rgb]{0.9,0.9,0.9}{
            \makecell[{{p{\linewidth}}}]{
                \texttt{\tiny{[GM$|$GM]}}
                \texttt{LIIGUTA: O, (4, 7)} \\
            }
        }
    }
    & & \\ \\

    \theutterance \stepcounter{utterance}  
    & & & \multicolumn{2}{p{0.3\linewidth}}{
        \cellcolor[rgb]{0.9,0.9,0.9}{
            \makecell[{{p{\linewidth}}}]{
                \texttt{\tiny{[GM$|$GM]}}
                \texttt{Karistus: (4,7) ei ole tühi, vaid sisaldab objekti 'M'.} \\
\\ 
\texttt{Teil on kollektiivselt kogunenud 6 karistust 16 karistusest.} \\
\texttt{Palun proovi uuesti!} \\
            }
        }
    }
    & & \\ \\

    \theutterance \stepcounter{utterance}  
    & & \multicolumn{4}{p{0.6\linewidth}}{
        \cellcolor[rgb]{0.9,0.9,0.9}{
            \makecell[{{p{\linewidth}}}]{
                \texttt{\tiny{[GM$\rangle$P2]}}
                \texttt{Karistus: (4,7) ei ole tühi, vaid sisaldab objekti 'M'.} \\
\\ 
\texttt{Teil on kollektiivselt kogunenud 6 karistust 16 karistusest.} \\
\texttt{Palun proovi uuesti!} \\
            }
        }
    }
    & \\ \\

    \theutterance \stepcounter{utterance}  
    & & & \multicolumn{4}{p{0.6\linewidth}}{
        \cellcolor[rgb]{0.9,0.9,0.9}{
            \makecell[{{p{\linewidth}}}]{
                \texttt{\tiny{[GM$\langle$P2]}}
                \texttt{LIIGUTA: O, (5, 7)} \\
            }
        }
    }
     \\ \\

    \theutterance \stepcounter{utterance}  
    & & & \multicolumn{2}{p{0.3\linewidth}}{
        \cellcolor[rgb]{0.9,0.9,0.9}{
            \makecell[{{p{\linewidth}}}]{
                \texttt{\tiny{[GM$|$GM]}}
                \texttt{LIIGUTA: O, (5, 7)} \\
            }
        }
    }
    & & \\ \\

    \theutterance \stepcounter{utterance}  
    & & & \multicolumn{2}{p{0.3\linewidth}}{
        \cellcolor[rgb]{0.9,0.9,0.9}{
            \makecell[{{p{\linewidth}}}]{
                \texttt{\tiny{[GM$|$GM]}}
                \texttt{Karistus: (5,7) ei ole tühi, vaid sisaldab objekti 'P'.} \\
\\ 
\texttt{Teil on kollektiivselt kogunenud 7 karistust 16 karistusest.} \\
\texttt{Palun proovi uuesti!} \\
            }
        }
    }
    & & \\ \\

    \theutterance \stepcounter{utterance}  
    & & \multicolumn{4}{p{0.6\linewidth}}{
        \cellcolor[rgb]{0.9,0.9,0.9}{
            \makecell[{{p{\linewidth}}}]{
                \texttt{\tiny{[GM$\rangle$P2]}}
                \texttt{Karistus: (5,7) ei ole tühi, vaid sisaldab objekti 'P'.} \\
\\ 
\texttt{Teil on kollektiivselt kogunenud 7 karistust 16 karistusest.} \\
\texttt{Palun proovi uuesti!} \\
            }
        }
    }
    & \\ \\

    \theutterance \stepcounter{utterance}  
    & & & \multicolumn{4}{p{0.6\linewidth}}{
        \cellcolor[rgb]{0.9,0.9,0.9}{
            \makecell[{{p{\linewidth}}}]{
                \texttt{\tiny{[GM$\langle$P2]}}
                \texttt{LIIGUTA: O, (6, 7)} \\
            }
        }
    }
     \\ \\

    \theutterance \stepcounter{utterance}  
    & & & \multicolumn{2}{p{0.3\linewidth}}{
        \cellcolor[rgb]{0.9,0.9,0.9}{
            \makecell[{{p{\linewidth}}}]{
                \texttt{\tiny{[GM$|$GM]}}
                \texttt{LIIGUTA: O, (6, 7)} \\
            }
        }
    }
    & & \\ \\

    \theutterance \stepcounter{utterance}  
    & & & \multicolumn{2}{p{0.3\linewidth}}{
        \cellcolor[rgb]{0.9,0.9,0.9}{
            \makecell[{{p{\linewidth}}}]{
                \texttt{\tiny{[GM$|$GM]}}
                \texttt{Karistus: (6,7) ei ole tühi, vaid sisaldab objekti '│'.} \\
\\ 
\texttt{Teil on kollektiivselt kogunenud 8 karistust 16 karistusest.} \\
\texttt{Palun proovi uuesti!} \\
            }
        }
    }
    & & \\ \\

    \theutterance \stepcounter{utterance}  
    & & \multicolumn{4}{p{0.6\linewidth}}{
        \cellcolor[rgb]{0.9,0.9,0.9}{
            \makecell[{{p{\linewidth}}}]{
                \texttt{\tiny{[GM$\rangle$P2]}}
                \texttt{Karistus: (6,7) ei ole tühi, vaid sisaldab objekti '│'.} \\
\\ 
\texttt{Teil on kollektiivselt kogunenud 8 karistust 16 karistusest.} \\
\texttt{Palun proovi uuesti!} \\
            }
        }
    }
    & \\ \\

    \theutterance \stepcounter{utterance}  
    & & & \multicolumn{4}{p{0.6\linewidth}}{
        \cellcolor[rgb]{0.9,0.9,0.9}{
            \makecell[{{p{\linewidth}}}]{
                \texttt{\tiny{[GM$\langle$P2]}}
                \texttt{LIIGUTA: O, (7, 7)} \\
            }
        }
    }
     \\ \\

    \theutterance \stepcounter{utterance}  
    & & & \multicolumn{2}{p{0.3\linewidth}}{
        \cellcolor[rgb]{0.9,0.9,0.9}{
            \makecell[{{p{\linewidth}}}]{
                \texttt{\tiny{[GM$|$GM]}}
                \texttt{LIIGUTA: O, (7, 7)} \\
            }
        }
    }
    & & \\ \\

    \theutterance \stepcounter{utterance}  
    & & & \multicolumn{2}{p{0.3\linewidth}}{
        \cellcolor[rgb]{0.9,0.9,0.9}{
            \makecell[{{p{\linewidth}}}]{
                \texttt{\tiny{[GM$|$GM]}}
                \texttt{O moved to (7, 7)} \\
            }
        }
    }
    & & \\ \\

    \theutterance \stepcounter{utterance}  
    & & \multicolumn{4}{p{0.6\linewidth}}{
        \cellcolor[rgb]{0.9,0.9,0.9}{
            \makecell[{{p{\linewidth}}}]{
                \texttt{\tiny{[P1$\langle$GM]}}
                \texttt{Edukalt teisaldatud 'A' asukohta (1,7). Teie uuendatud ruudustik näeb välja selline:} \\
\texttt{```} \\
\texttt{ 1234567} \\
\texttt{╔═╤═════╗} \\
\texttt{║◌│◌◌◌◌◌║ 1} \\
\texttt{╟─┤◌T◌◌◌║ 2} \\
\texttt{║O└─┐◌◌◌║ 3} \\
\texttt{║◌◌◌├─┐◌║ 4} \\
\texttt{║◌◌S│◌│◌║ 5} \\
\texttt{╟─┬─┴─┼─╢ 6} \\
\texttt{║A│◌DP│M║ 7} \\
\texttt{╚═╧═══╧═╝} \\
\texttt{```} \\
\texttt{Sa mängid praegu vooru 8 maksimaalsest 28 voorust.} \\
\texttt{Teil on kollektiivselt kogunenud 8 karistust 16 karistusest.} \\
\texttt{Teine mängija liigutas objekti oma ruudustikul.} \\
\\ 
\\ 
\texttt{Mis on sinu järgmine käsk?} \\
            }
        }
    }
    & \\ \\

    \theutterance \stepcounter{utterance}  
    & \multicolumn{4}{p{0.6\linewidth}}{
        \cellcolor[rgb]{0.9,0.9,0.9}{
            \makecell[{{p{\linewidth}}}]{
                \texttt{\tiny{[P1$\rangle$GM]}}
                \texttt{LIIGUTA: O, (4, 7)} \\
            }
        }
    }
    & & \\ \\

    \theutterance \stepcounter{utterance}  
    & & & \multicolumn{2}{p{0.3\linewidth}}{
        \cellcolor[rgb]{0.9,0.9,0.9}{
            \makecell[{{p{\linewidth}}}]{
                \texttt{\tiny{[GM$|$GM]}}
                \texttt{LIIGUTA: O, (4, 7)} \\
            }
        }
    }
    & & \\ \\

    \theutterance \stepcounter{utterance}  
    & & & \multicolumn{2}{p{0.3\linewidth}}{
        \cellcolor[rgb]{0.9,0.9,0.9}{
            \makecell[{{p{\linewidth}}}]{
                \texttt{\tiny{[GM$|$GM]}}
                \texttt{Karistus: (4,7) ei ole tühi, vaid sisaldab objekti 'D'.} \\
\\ 
\texttt{Teil on kollektiivselt kogunenud 9 karistust 16 karistusest.} \\
\texttt{Palun proovi uuesti!} \\
            }
        }
    }
    & & \\ \\

    \theutterance \stepcounter{utterance}  
    & & \multicolumn{4}{p{0.6\linewidth}}{
        \cellcolor[rgb]{0.9,0.9,0.9}{
            \makecell[{{p{\linewidth}}}]{
                \texttt{\tiny{[P1$\langle$GM]}}
                \texttt{Karistus: (4,7) ei ole tühi, vaid sisaldab objekti 'D'.} \\
\\ 
\texttt{Teil on kollektiivselt kogunenud 9 karistust 16 karistusest.} \\
\texttt{Palun proovi uuesti!} \\
            }
        }
    }
    & \\ \\

    \theutterance \stepcounter{utterance}  
    & \multicolumn{4}{p{0.6\linewidth}}{
        \cellcolor[rgb]{0.9,0.9,0.9}{
            \makecell[{{p{\linewidth}}}]{
                \texttt{\tiny{[P1$\rangle$GM]}}
                \texttt{LIIGUTA: O, (3, 7)} \\
            }
        }
    }
    & & \\ \\

    \theutterance \stepcounter{utterance}  
    & & & \multicolumn{2}{p{0.3\linewidth}}{
        \cellcolor[rgb]{0.9,0.9,0.9}{
            \makecell[{{p{\linewidth}}}]{
                \texttt{\tiny{[GM$|$GM]}}
                \texttt{LIIGUTA: O, (3, 7)} \\
            }
        }
    }
    & & \\ \\

    \theutterance \stepcounter{utterance}  
    & & & \multicolumn{2}{p{0.3\linewidth}}{
        \cellcolor[rgb]{0.9,0.9,0.9}{
            \makecell[{{p{\linewidth}}}]{
                \texttt{\tiny{[GM$|$GM]}}
                \texttt{O moved to (3, 7)} \\
            }
        }
    }
    & & \\ \\

    \theutterance \stepcounter{utterance}  
    & & \multicolumn{4}{p{0.6\linewidth}}{
        \cellcolor[rgb]{0.9,0.9,0.9}{
            \makecell[{{p{\linewidth}}}]{
                \texttt{\tiny{[GM$\rangle$P2]}}
                \texttt{Edukalt teisaldatud 'O' asukohta (7,7). Teie uuendatud ruudustik näeb välja selline:} \\
\texttt{```} \\
\texttt{ 1234567} \\
\texttt{╔═╤═════╗} \\
\texttt{║◌│S◌◌◌◌║ 1} \\
\texttt{╟─┤◌◌◌◌◌║ 2} \\
\texttt{║◌└─┐◌◌◌║ 3} \\
\texttt{║◌◌◌├─┐◌║ 4} \\
\texttt{║◌T◌│◌│◌║ 5} \\
\texttt{╟─┬─┴─┼─╢ 6} \\
\texttt{║A│DMP│O║ 7} \\
\texttt{╚═╧═══╧═╝} \\
\texttt{```} \\
\texttt{Sa mängid praegu vooru 9 maksimaalsest 28 voorust.} \\
\texttt{Teil on kollektiivselt kogunenud 9 karistust 16 karistusest.} \\
\texttt{Teine mängija liigutas objekti oma ruudustikul.} \\
\\ 
\\ 
\texttt{Mis on sinu järgmine käsk?} \\
            }
        }
    }
    & \\ \\

    \theutterance \stepcounter{utterance}  
    & & & \multicolumn{4}{p{0.6\linewidth}}{
        \cellcolor[rgb]{0.9,0.9,0.9}{
            \makecell[{{p{\linewidth}}}]{
                \texttt{\tiny{[GM$\langle$P2]}}
                \texttt{LIIGUTA: S, (6, 7)} \\
            }
        }
    }
     \\ \\

    \theutterance \stepcounter{utterance}  
    & & & \multicolumn{2}{p{0.3\linewidth}}{
        \cellcolor[rgb]{0.9,0.9,0.9}{
            \makecell[{{p{\linewidth}}}]{
                \texttt{\tiny{[GM$|$GM]}}
                \texttt{LIIGUTA: S, (6, 7)} \\
            }
        }
    }
    & & \\ \\

    \theutterance \stepcounter{utterance}  
    & & & \multicolumn{2}{p{0.3\linewidth}}{
        \cellcolor[rgb]{0.9,0.9,0.9}{
            \makecell[{{p{\linewidth}}}]{
                \texttt{\tiny{[GM$|$GM]}}
                \texttt{Karistus: (6,7) ei ole tühi, vaid sisaldab objekti '│'.} \\
\\ 
\texttt{Teil on kollektiivselt kogunenud 10 karistust 16 karistusest.} \\
\texttt{Palun proovi uuesti!} \\
            }
        }
    }
    & & \\ \\

    \theutterance \stepcounter{utterance}  
    & & \multicolumn{4}{p{0.6\linewidth}}{
        \cellcolor[rgb]{0.9,0.9,0.9}{
            \makecell[{{p{\linewidth}}}]{
                \texttt{\tiny{[GM$\rangle$P2]}}
                \texttt{Karistus: (6,7) ei ole tühi, vaid sisaldab objekti '│'.} \\
\\ 
\texttt{Teil on kollektiivselt kogunenud 10 karistust 16 karistusest.} \\
\texttt{Palun proovi uuesti!} \\
            }
        }
    }
    & \\ \\

    \theutterance \stepcounter{utterance}  
    & & & \multicolumn{4}{p{0.6\linewidth}}{
        \cellcolor[rgb]{0.9,0.9,0.9}{
            \makecell[{{p{\linewidth}}}]{
                \texttt{\tiny{[GM$\langle$P2]}}
                \texttt{ÜTLE: Minu rida 7 veerg 6 on seina osa ('│'), seega ei saa seal objekti paigutada. Tegin viga plaanis – S peaks minema (6,7), aga see on blokeeritud. Muudan S sihtkohta (1,7), liigutan A (1,7){-}>(2,7) vabastamiseks. Liigutan A (1,7) {-}> (2,7).} \\
            }
        }
    }
     \\ \\

    \theutterance \stepcounter{utterance}  
    & & & \multicolumn{2}{p{0.3\linewidth}}{
        \cellcolor[rgb]{0.9,0.9,0.9}{
            \makecell[{{p{\linewidth}}}]{
                \texttt{\tiny{[GM$|$GM]}}
                \texttt{ÜTLE: Minu rida 7 veerg 6 on seina osa ('│'), seega ei saa seal objekti paigutada. Tegin viga plaanis – S peaks minema (6,7), aga see on blokeeritud. Muudan S sihtkohta (1,7), liigutan A (1,7){-}>(2,7) vabastamiseks. Liigutan A (1,7) {-}> (2,7).} \\
            }
        }
    }
    & & \\ \\

    \theutterance \stepcounter{utterance}  
    & & \multicolumn{4}{p{0.6\linewidth}}{
        \cellcolor[rgb]{0.9,0.9,0.9}{
            \makecell[{{p{\linewidth}}}]{
                \texttt{\tiny{[P1$\langle$GM]}}
                \texttt{Edukalt teisaldatud 'O' asukohta (3,7). Teie uuendatud ruudustik näeb välja selline:} \\
\texttt{```} \\
\texttt{ 1234567} \\
\texttt{╔═╤═════╗} \\
\texttt{║◌│◌◌◌◌◌║ 1} \\
\texttt{╟─┤◌T◌◌◌║ 2} \\
\texttt{║◌└─┐◌◌◌║ 3} \\
\texttt{║◌◌◌├─┐◌║ 4} \\
\texttt{║◌◌S│◌│◌║ 5} \\
\texttt{╟─┬─┴─┼─╢ 6} \\
\texttt{║A│ODP│M║ 7} \\
\texttt{╚═╧═══╧═╝} \\
\texttt{```} \\
\texttt{Sa mängid praegu vooru 9 maksimaalsest 28 voorust.} \\
\texttt{Teil on kollektiivselt kogunenud 10 karistust 16 karistusest.} \\
\texttt{Teine mängija saatis selle sõnumi:} \\
\\ 
\\ 
\texttt{"Minu rida 7 veerg 6 on seina osa ('│'), seega ei saa seal objekti paigutada. Tegin viga plaanis – S peaks minema (6,7), aga see on blokeeritud. Muudan S sihtkohta (1,7), liigutan A (1,7){-}>(2,7) vabastamiseks. Liigutan A (1,7) {-}> (2,7)."} \\
\\ 
\\ 
\texttt{Mis on sinu järgmine käsk?} \\
            }
        }
    }
    & \\ \\

    \theutterance \stepcounter{utterance}  
    & \multicolumn{4}{p{0.6\linewidth}}{
        \cellcolor[rgb]{0.9,0.9,0.9}{
            \makecell[{{p{\linewidth}}}]{
                \texttt{\tiny{[P1$\rangle$GM]}}
                \texttt{LIIGUTA: D, (2, 7)} \\
            }
        }
    }
    & & \\ \\

    \theutterance \stepcounter{utterance}  
    & & & \multicolumn{2}{p{0.3\linewidth}}{
        \cellcolor[rgb]{0.9,0.9,0.9}{
            \makecell[{{p{\linewidth}}}]{
                \texttt{\tiny{[GM$|$GM]}}
                \texttt{LIIGUTA: D, (2, 7)} \\
            }
        }
    }
    & & \\ \\

    \theutterance \stepcounter{utterance}  
    & & & \multicolumn{2}{p{0.3\linewidth}}{
        \cellcolor[rgb]{0.9,0.9,0.9}{
            \makecell[{{p{\linewidth}}}]{
                \texttt{\tiny{[GM$|$GM]}}
                \texttt{Karistus: (2,7) ei ole tühi, vaid sisaldab objekti '│'.} \\
\\ 
\texttt{Teil on kollektiivselt kogunenud 11 karistust 16 karistusest.} \\
\texttt{Palun proovi uuesti!} \\
            }
        }
    }
    & & \\ \\

    \theutterance \stepcounter{utterance}  
    & & \multicolumn{4}{p{0.6\linewidth}}{
        \cellcolor[rgb]{0.9,0.9,0.9}{
            \makecell[{{p{\linewidth}}}]{
                \texttt{\tiny{[P1$\langle$GM]}}
                \texttt{Karistus: (2,7) ei ole tühi, vaid sisaldab objekti '│'.} \\
\\ 
\texttt{Teil on kollektiivselt kogunenud 11 karistust 16 karistusest.} \\
\texttt{Palun proovi uuesti!} \\
            }
        }
    }
    & \\ \\

    \theutterance \stepcounter{utterance}  
    & \multicolumn{4}{p{0.6\linewidth}}{
        \cellcolor[rgb]{0.9,0.9,0.9}{
            \makecell[{{p{\linewidth}}}]{
                \texttt{\tiny{[P1$\rangle$GM]}}
                \texttt{LIIGUTA: D, (6, 7)} \\
            }
        }
    }
    & & \\ \\

    \theutterance \stepcounter{utterance}  
    & & & \multicolumn{2}{p{0.3\linewidth}}{
        \cellcolor[rgb]{0.9,0.9,0.9}{
            \makecell[{{p{\linewidth}}}]{
                \texttt{\tiny{[GM$|$GM]}}
                \texttt{LIIGUTA: D, (6, 7)} \\
            }
        }
    }
    & & \\ \\

    \theutterance \stepcounter{utterance}  
    & & & \multicolumn{2}{p{0.3\linewidth}}{
        \cellcolor[rgb]{0.9,0.9,0.9}{
            \makecell[{{p{\linewidth}}}]{
                \texttt{\tiny{[GM$|$GM]}}
                \texttt{Karistus: (6,7) ei ole tühi, vaid sisaldab objekti '│'.} \\
\\ 
\texttt{Teil on kollektiivselt kogunenud 12 karistust 16 karistusest.} \\
\texttt{Palun proovi uuesti!} \\
            }
        }
    }
    & & \\ \\

    \theutterance \stepcounter{utterance}  
    & & \multicolumn{4}{p{0.6\linewidth}}{
        \cellcolor[rgb]{0.9,0.9,0.9}{
            \makecell[{{p{\linewidth}}}]{
                \texttt{\tiny{[P1$\langle$GM]}}
                \texttt{Karistus: (6,7) ei ole tühi, vaid sisaldab objekti '│'.} \\
\\ 
\texttt{Teil on kollektiivselt kogunenud 12 karistust 16 karistusest.} \\
\texttt{Palun proovi uuesti!} \\
            }
        }
    }
    & \\ \\

    \theutterance \stepcounter{utterance}  
    & \multicolumn{4}{p{0.6\linewidth}}{
        \cellcolor[rgb]{0.9,0.9,0.9}{
            \makecell[{{p{\linewidth}}}]{
                \texttt{\tiny{[P1$\rangle$GM]}}
                \texttt{ÜTLE: Minu rida 7 on blokeeritud: veergud 2, 4, 6 on seinaelemendid. Ainult veerud 1, 3, 5, 7 on vabad. Kuna meil on 7 objekti, ei mahu need rea 7. Teie S on ka blokeeritud. Peaksime valima teise rea. Rida 4 näeb vaba välja (veergud 1{-}7). Kas sobib, kui paigutame kõik read 4? Sihtkoordinaadid: A(1,4), D(2,4), M(3,4), O(4,4), P(5,4), S(6,4), T(7,4).} \\
            }
        }
    }
    & & \\ \\

    \theutterance \stepcounter{utterance}  
    & & & \multicolumn{2}{p{0.3\linewidth}}{
        \cellcolor[rgb]{0.9,0.9,0.9}{
            \makecell[{{p{\linewidth}}}]{
                \texttt{\tiny{[GM$|$GM]}}
                \texttt{ÜTLE: Minu rida 7 on blokeeritud: veergud 2, 4, 6 on seinaelemendid. Ainult veerud 1, 3, 5, 7 on vabad. Kuna meil on 7 objekti, ei mahu need rea 7. Teie S on ka blokeeritud. Peaksime valima teise rea. Rida 4 näeb vaba välja (veergud 1{-}7). Kas sobib, kui paigutame kõik read 4? Sihtkoordinaadid: A(1,4), D(2,4), M(3,4), O(4,4), P(5,4), S(6,4), T(7,4).} \\
            }
        }
    }
    & & \\ \\

    \theutterance \stepcounter{utterance}  
    & & \multicolumn{4}{p{0.6\linewidth}}{
        \cellcolor[rgb]{0.9,0.9,0.9}{
            \makecell[{{p{\linewidth}}}]{
                \texttt{\tiny{[GM$\rangle$P2]}}
                \texttt{Sinu sõnum on teisele mängijale edastatud.} \\
\\ 
\\ 
\texttt{Sa mängid praegu vooru 10 maksimaalsest 28 voorust.} \\
\texttt{Teil on kollektiivselt kogunenud 12 karistust 16 karistusest.} \\
\texttt{Teine mängija saatis selle sõnumi:} \\
\\ 
\\ 
\texttt{"Minu rida 7 on blokeeritud: veergud 2, 4, 6 on seinaelemendid. Ainult veerud 1, 3, 5, 7 on vabad. Kuna meil on 7 objekti, ei mahu need rea 7. Teie S on ka blokeeritud. Peaksime valima teise rea. Rida 4 näeb vaba välja (veergud 1{-}7). Kas sobib, kui paigutame kõik read 4? Sihtkoordinaadid: A(1,4), D(2,4), M(3,4), O(4,4), P(5,4), S(6,4), T(7,4)."} \\
\\ 
\\ 
\texttt{Mis on sinu järgmine käsk?} \\
            }
        }
    }
    & \\ \\

    \theutterance \stepcounter{utterance}  
    & & & \multicolumn{4}{p{0.6\linewidth}}{
        \cellcolor[rgb]{0.9,0.9,0.9}{
            \makecell[{{p{\linewidth}}}]{
                \texttt{\tiny{[GM$\langle$P2]}}
                \texttt{ÜTLE: Rida 4 on minu laual osaliselt hõivatud: veerg 5 on '├' (sein), veerg 7 on '◌' (vaba). Seega ei ole kogu rida 4 vaba. Rida 1 on minu laual täiesti vaba (veergud 1{-}7 on kõik '◌'). Soovitan paigutada kõik read 1. Sihtkoordinaadid: A(1,1), D(2,1), M(3,1), O(4,1), P(5,1), S(6,1), T(7,1). Kas see sobib?} \\
            }
        }
    }
     \\ \\

    \theutterance \stepcounter{utterance}  
    & & & \multicolumn{2}{p{0.3\linewidth}}{
        \cellcolor[rgb]{0.9,0.9,0.9}{
            \makecell[{{p{\linewidth}}}]{
                \texttt{\tiny{[GM$|$GM]}}
                \texttt{ÜTLE: Rida 4 on minu laual osaliselt hõivatud: veerg 5 on '├' (sein), veerg 7 on '◌' (vaba). Seega ei ole kogu rida 4 vaba. Rida 1 on minu laual täiesti vaba (veergud 1{-}7 on kõik '◌'). Soovitan paigutada kõik read 1. Sihtkoordinaadid: A(1,1), D(2,1), M(3,1), O(4,1), P(5,1), S(6,1), T(7,1). Kas see sobib?} \\
            }
        }
    }
    & & \\ \\

    \theutterance \stepcounter{utterance}  
    & & \multicolumn{4}{p{0.6\linewidth}}{
        \cellcolor[rgb]{0.9,0.9,0.9}{
            \makecell[{{p{\linewidth}}}]{
                \texttt{\tiny{[P1$\langle$GM]}}
                \texttt{Sinu sõnum on teisele mängijale edastatud.} \\
\\ 
\\ 
\texttt{Sa mängid praegu vooru 10 maksimaalsest 28 voorust.} \\
\texttt{Teil on kollektiivselt kogunenud 12 karistust 16 karistusest.} \\
\texttt{Teine mängija saatis selle sõnumi:} \\
\\ 
\\ 
\texttt{"Rida 4 on minu laual osaliselt hõivatud: veerg 5 on '├' (sein), veerg 7 on '◌' (vaba). Seega ei ole kogu rida 4 vaba. Rida 1 on minu laual täiesti vaba (veergud 1{-}7 on kõik '◌'). Soovitan paigutada kõik read 1. Sihtkoordinaadid: A(1,1), D(2,1), M(3,1), O(4,1), P(5,1), S(6,1), T(7,1). Kas see sobib?"} \\
\\ 
\\ 
\texttt{Mis on sinu järgmine käsk?} \\
            }
        }
    }
    & \\ \\

    \theutterance \stepcounter{utterance}  
    & \multicolumn{4}{p{0.6\linewidth}}{
        \cellcolor[rgb]{0.9,0.9,0.9}{
            \makecell[{{p{\linewidth}}}]{
                \texttt{\tiny{[P1$\rangle$GM]}}
                \texttt{ÜTLE: Rida 1 on minu laual hõivatud (veerg 2='│', 3='A'). Rida 5 on minu laual täiesti vaba (veergud 1{-}7 kõik '◌'). Kas teie rida 5 on vaba? Sihtkoordinaadid: A(1,5), D(2,5), M(3,5), O(4,5), P(5,5), S(6,5), T(7,5).} \\
            }
        }
    }
    & & \\ \\

    \theutterance \stepcounter{utterance}  
    & & & \multicolumn{2}{p{0.3\linewidth}}{
        \cellcolor[rgb]{0.9,0.9,0.9}{
            \makecell[{{p{\linewidth}}}]{
                \texttt{\tiny{[GM$|$GM]}}
                \texttt{ÜTLE: Rida 1 on minu laual hõivatud (veerg 2='│', 3='A'). Rida 5 on minu laual täiesti vaba (veergud 1{-}7 kõik '◌'). Kas teie rida 5 on vaba? Sihtkoordinaadid: A(1,5), D(2,5), M(3,5), O(4,5), P(5,5), S(6,5), T(7,5).} \\
            }
        }
    }
    & & \\ \\

    \theutterance \stepcounter{utterance}  
    & & \multicolumn{4}{p{0.6\linewidth}}{
        \cellcolor[rgb]{0.9,0.9,0.9}{
            \makecell[{{p{\linewidth}}}]{
                \texttt{\tiny{[GM$\rangle$P2]}}
                \texttt{Sinu sõnum on teisele mängijale edastatud.} \\
\\ 
\\ 
\texttt{Sa mängid praegu vooru 11 maksimaalsest 28 voorust.} \\
\texttt{Teil on kollektiivselt kogunenud 12 karistust 16 karistusest.} \\
\texttt{Teine mängija saatis selle sõnumi:} \\
\\ 
\\ 
\texttt{"Rida 1 on minu laual hõivatud (veerg 2='│', 3='A'). Rida 5 on minu laual täiesti vaba (veergud 1{-}7 kõik '◌'). Kas teie rida 5 on vaba? Sihtkoordinaadid: A(1,5), D(2,5), M(3,5), O(4,5), P(5,5), S(6,5), T(7,5)."} \\
\\ 
\\ 
\texttt{Mis on sinu järgmine käsk?} \\
            }
        }
    }
    & \\ \\

    \theutterance \stepcounter{utterance}  
    & & & \multicolumn{4}{p{0.6\linewidth}}{
        \cellcolor[rgb]{0.9,0.9,0.9}{
            \makecell[{{p{\linewidth}}}]{
                \texttt{\tiny{[GM$\langle$P2]}}
                \texttt{ÜTLE: Minu rida 5 on osaliselt hõivatud: veerg 2 on 'T', veerg 4 on '│', veerg 6 on '│', veerg 7 on 'O'. Veerud 1, 3, 5 on vabad. Rida 3 on minu laual täiesti vaba (veergud 1{-}7 kõik '◌'). Soovitan rida 3. Sihtkoordinaadid: A(1,3), D(2,3), M(3,3), O(4,3), P(5,3), S(6,3), T(7,3). Kas see sobib teile?} \\
            }
        }
    }
     \\ \\

    \theutterance \stepcounter{utterance}  
    & & & \multicolumn{2}{p{0.3\linewidth}}{
        \cellcolor[rgb]{0.9,0.9,0.9}{
            \makecell[{{p{\linewidth}}}]{
                \texttt{\tiny{[GM$|$GM]}}
                \texttt{ÜTLE: Minu rida 5 on osaliselt hõivatud: veerg 2 on 'T', veerg 4 on '│', veerg 6 on '│', veerg 7 on 'O'. Veerud 1, 3, 5 on vabad. Rida 3 on minu laual täiesti vaba (veergud 1{-}7 kõik '◌'). Soovitan rida 3. Sihtkoordinaadid: A(1,3), D(2,3), M(3,3), O(4,3), P(5,3), S(6,3), T(7,3). Kas see sobib teile?} \\
            }
        }
    }
    & & \\ \\

    \theutterance \stepcounter{utterance}  
    & & \multicolumn{4}{p{0.6\linewidth}}{
        \cellcolor[rgb]{0.9,0.9,0.9}{
            \makecell[{{p{\linewidth}}}]{
                \texttt{\tiny{[P1$\langle$GM]}}
                \texttt{Sinu sõnum on teisele mängijale edastatud.} \\
\\ 
\\ 
\texttt{Sa mängid praegu vooru 11 maksimaalsest 28 voorust.} \\
\texttt{Teil on kollektiivselt kogunenud 12 karistust 16 karistusest.} \\
\texttt{Teine mängija saatis selle sõnumi:} \\
\\ 
\\ 
\texttt{"Minu rida 5 on osaliselt hõivatud: veerg 2 on 'T', veerg 4 on '│', veerg 6 on '│', veerg 7 on 'O'. Veerud 1, 3, 5 on vabad. Rida 3 on minu laual täiesti vaba (veergud 1{-}7 kõik '◌'). Soovitan rida 3. Sihtkoordinaadid: A(1,3), D(2,3), M(3,3), O(4,3), P(5,3), S(6,3), T(7,3). Kas see sobib teile?"} \\
\\ 
\\ 
\texttt{Mis on sinu järgmine käsk?} \\
            }
        }
    }
    & \\ \\

    \theutterance \stepcounter{utterance}  
    & \multicolumn{4}{p{0.6\linewidth}}{
        \cellcolor[rgb]{0.9,0.9,0.9}{
            \makecell[{{p{\linewidth}}}]{
                \texttt{\tiny{[P1$\rangle$GM]}}
                \texttt{ÜTLE: Minu rida 3 on hõivatud (veerg 1='O', veerg 4='┐'). Minu vabad read on 1, 4, 5. Teie vabad read on 1, 3. Ainult rida 1 on ühine vaba rida, aga minu veerg 2 on seina. Lahendus: paigutame read 1 veergudesse 1,3,4,5,6,7 (6 objekti) ja ühe objekti (nt T) veergu 2 rea 4 (mis on mõlemal vaba). Kas see sobib? Siht: A(1,1), D(3,1), M(4,1), O(5,1), P(6,1), S(7,1), T(2,4).} \\
            }
        }
    }
    & & \\ \\

    \theutterance \stepcounter{utterance}  
    & & & \multicolumn{2}{p{0.3\linewidth}}{
        \cellcolor[rgb]{0.9,0.9,0.9}{
            \makecell[{{p{\linewidth}}}]{
                \texttt{\tiny{[GM$|$GM]}}
                \texttt{ÜTLE: Minu rida 3 on hõivatud (veerg 1='O', veerg 4='┐'). Minu vabad read on 1, 4, 5. Teie vabad read on 1, 3. Ainult rida 1 on ühine vaba rida, aga minu veerg 2 on seina. Lahendus: paigutame read 1 veergudesse 1,3,4,5,6,7 (6 objekti) ja ühe objekti (nt T) veergu 2 rea 4 (mis on mõlemal vaba). Kas see sobib? Siht: A(1,1), D(3,1), M(4,1), O(5,1), P(6,1), S(7,1), T(2,4).} \\
            }
        }
    }
    & & \\ \\

    \theutterance \stepcounter{utterance}  
    & & \multicolumn{4}{p{0.6\linewidth}}{
        \cellcolor[rgb]{0.9,0.9,0.9}{
            \makecell[{{p{\linewidth}}}]{
                \texttt{\tiny{[GM$\rangle$P2]}}
                \texttt{Sinu sõnum on teisele mängijale edastatud.} \\
\\ 
\\ 
\texttt{Sa mängid praegu vooru 12 maksimaalsest 28 voorust.} \\
\texttt{Teil on kollektiivselt kogunenud 12 karistust 16 karistusest.} \\
\texttt{Teine mängija saatis selle sõnumi:} \\
\\ 
\\ 
\texttt{"Minu rida 3 on hõivatud (veerg 1='O', veerg 4='┐'). Minu vabad read on 1, 4, 5. Teie vabad read on 1, 3. Ainult rida 1 on ühine vaba rida, aga minu veerg 2 on seina. Lahendus: paigutame read 1 veergudesse 1,3,4,5,6,7 (6 objekti) ja ühe objekti (nt T) veergu 2 rea 4 (mis on mõlemal vaba). Kas see sobib? Siht: A(1,1), D(3,1), M(4,1), O(5,1), P(6,1), S(7,1), T(2,4)."} \\
\\ 
\\ 
\texttt{Mis on sinu järgmine käsk?} \\
            }
        }
    }
    & \\ \\

    \theutterance \stepcounter{utterance}  
    & & & \multicolumn{4}{p{0.6\linewidth}}{
        \cellcolor[rgb]{0.9,0.9,0.9}{
            \makecell[{{p{\linewidth}}}]{
                \texttt{\tiny{[GM$\langle$P2]}}
                \texttt{ÜTLE: Minu rida 1 on täiesti vaba (veerud 1{-}7). Teie rida 1 veerg 2 on seina. Teie pakkumine (6 objekti rida 1 + T(2,4)) on hea kompromiss. Kinnitan: sihtkoordinaadid A(1,1), D(3,1), M(4,1), O(5,1), P(6,1), S(7,1), T(2,4). Alustan kohe A liigutusega (6,3){-}>(1,1).} \\
            }
        }
    }
     \\ \\

    \theutterance \stepcounter{utterance}  
    & & & \multicolumn{2}{p{0.3\linewidth}}{
        \cellcolor[rgb]{0.9,0.9,0.9}{
            \makecell[{{p{\linewidth}}}]{
                \texttt{\tiny{[GM$|$GM]}}
                \texttt{ÜTLE: Minu rida 1 on täiesti vaba (veerud 1{-}7). Teie rida 1 veerg 2 on seina. Teie pakkumine (6 objekti rida 1 + T(2,4)) on hea kompromiss. Kinnitan: sihtkoordinaadid A(1,1), D(3,1), M(4,1), O(5,1), P(6,1), S(7,1), T(2,4). Alustan kohe A liigutusega (6,3){-}>(1,1).} \\
            }
        }
    }
    & & \\ \\

    \theutterance \stepcounter{utterance}  
    & & \multicolumn{4}{p{0.6\linewidth}}{
        \cellcolor[rgb]{0.9,0.9,0.9}{
            \makecell[{{p{\linewidth}}}]{
                \texttt{\tiny{[P1$\langle$GM]}}
                \texttt{Sinu sõnum on teisele mängijale edastatud.} \\
\\ 
\\ 
\texttt{Sa mängid praegu vooru 12 maksimaalsest 28 voorust.} \\
\texttt{Teil on kollektiivselt kogunenud 12 karistust 16 karistusest.} \\
\texttt{Teine mängija saatis selle sõnumi:} \\
\\ 
\\ 
\texttt{"Minu rida 1 on täiesti vaba (veerud 1{-}7). Teie rida 1 veerg 2 on seina. Teie pakkumine (6 objekti rida 1 + T(2,4)) on hea kompromiss. Kinnitan: sihtkoordinaadid A(1,1), D(3,1), M(4,1), O(5,1), P(6,1), S(7,1), T(2,4). Alustan kohe A liigutusega (6,3){-}>(1,1)."} \\
\\ 
\\ 
\texttt{Mis on sinu järgmine käsk?} \\
            }
        }
    }
    & \\ \\

    \theutterance \stepcounter{utterance}  
    & \multicolumn{4}{p{0.6\linewidth}}{
        \cellcolor[rgb]{0.9,0.9,0.9}{
            \makecell[{{p{\linewidth}}}]{
                \texttt{\tiny{[P1$\rangle$GM]}}
                \texttt{LIIGUTA: A, (1, 1)} \\
            }
        }
    }
    & & \\ \\

    \theutterance \stepcounter{utterance}  
    & & & \multicolumn{2}{p{0.3\linewidth}}{
        \cellcolor[rgb]{0.9,0.9,0.9}{
            \makecell[{{p{\linewidth}}}]{
                \texttt{\tiny{[GM$|$GM]}}
                \texttt{LIIGUTA: A, (1, 1)} \\
            }
        }
    }
    & & \\ \\

    \theutterance \stepcounter{utterance}  
    & & & \multicolumn{2}{p{0.3\linewidth}}{
        \cellcolor[rgb]{0.9,0.9,0.9}{
            \makecell[{{p{\linewidth}}}]{
                \texttt{\tiny{[GM$|$GM]}}
                \texttt{A moved to (1, 1)} \\
            }
        }
    }
    & & \\ \\

    \theutterance \stepcounter{utterance}  
    & & \multicolumn{4}{p{0.6\linewidth}}{
        \cellcolor[rgb]{0.9,0.9,0.9}{
            \makecell[{{p{\linewidth}}}]{
                \texttt{\tiny{[GM$\rangle$P2]}}
                \texttt{Sinu sõnum on teisele mängijale edastatud.} \\
\\ 
\\ 
\texttt{Sa mängid praegu vooru 13 maksimaalsest 28 voorust.} \\
\texttt{Teil on kollektiivselt kogunenud 12 karistust 16 karistusest.} \\
\texttt{Teine mängija liigutas objekti oma ruudustikul.} \\
\\ 
\\ 
\texttt{Mis on sinu järgmine käsk?} \\
            }
        }
    }
    & \\ \\

    \theutterance \stepcounter{utterance}  
    & & & \multicolumn{4}{p{0.6\linewidth}}{
        \cellcolor[rgb]{0.9,0.9,0.9}{
            \makecell[{{p{\linewidth}}}]{
                \texttt{\tiny{[GM$\langle$P2]}}
                \texttt{LIIGUTA: A, (1, 1)} \\
            }
        }
    }
     \\ \\

    \theutterance \stepcounter{utterance}  
    & & & \multicolumn{2}{p{0.3\linewidth}}{
        \cellcolor[rgb]{0.9,0.9,0.9}{
            \makecell[{{p{\linewidth}}}]{
                \texttt{\tiny{[GM$|$GM]}}
                \texttt{LIIGUTA: A, (1, 1)} \\
            }
        }
    }
    & & \\ \\

    \theutterance \stepcounter{utterance}  
    & & & \multicolumn{2}{p{0.3\linewidth}}{
        \cellcolor[rgb]{0.9,0.9,0.9}{
            \makecell[{{p{\linewidth}}}]{
                \texttt{\tiny{[GM$|$GM]}}
                \texttt{A moved to (1, 1)} \\
            }
        }
    }
    & & \\ \\

    \theutterance \stepcounter{utterance}  
    & & \multicolumn{4}{p{0.6\linewidth}}{
        \cellcolor[rgb]{0.9,0.9,0.9}{
            \makecell[{{p{\linewidth}}}]{
                \texttt{\tiny{[P1$\langle$GM]}}
                \texttt{Edukalt teisaldatud 'A' asukohta (1,1). Teie uuendatud ruudustik näeb välja selline:} \\
\texttt{```} \\
\texttt{ 1234567} \\
\texttt{╔═╤═════╗} \\
\texttt{║A│◌◌◌◌◌║ 1} \\
\texttt{╟─┤◌T◌◌◌║ 2} \\
\texttt{║◌└─┐◌◌◌║ 3} \\
\texttt{║◌◌◌├─┐◌║ 4} \\
\texttt{║◌◌S│◌│◌║ 5} \\
\texttt{╟─┬─┴─┼─╢ 6} \\
\texttt{║◌│ODP│M║ 7} \\
\texttt{╚═╧═══╧═╝} \\
\texttt{```} \\
\texttt{Sa mängid praegu vooru 13 maksimaalsest 28 voorust.} \\
\texttt{Teil on kollektiivselt kogunenud 12 karistust 16 karistusest.} \\
\texttt{Teine mängija liigutas objekti oma ruudustikul.} \\
\\ 
\\ 
\texttt{Mis on sinu järgmine käsk?} \\
            }
        }
    }
    & \\ \\

    \theutterance \stepcounter{utterance}  
    & \multicolumn{4}{p{0.6\linewidth}}{
        \cellcolor[rgb]{0.9,0.9,0.9}{
            \makecell[{{p{\linewidth}}}]{
                \texttt{\tiny{[P1$\rangle$GM]}}
                \texttt{LIIGUTA: D, (3, 1)} \\
            }
        }
    }
    & & \\ \\

    \theutterance \stepcounter{utterance}  
    & & & \multicolumn{2}{p{0.3\linewidth}}{
        \cellcolor[rgb]{0.9,0.9,0.9}{
            \makecell[{{p{\linewidth}}}]{
                \texttt{\tiny{[GM$|$GM]}}
                \texttt{LIIGUTA: D, (3, 1)} \\
            }
        }
    }
    & & \\ \\

    \theutterance \stepcounter{utterance}  
    & & & \multicolumn{2}{p{0.3\linewidth}}{
        \cellcolor[rgb]{0.9,0.9,0.9}{
            \makecell[{{p{\linewidth}}}]{
                \texttt{\tiny{[GM$|$GM]}}
                \texttt{D moved to (3, 1)} \\
            }
        }
    }
    & & \\ \\

    \theutterance \stepcounter{utterance}  
    & & \multicolumn{4}{p{0.6\linewidth}}{
        \cellcolor[rgb]{0.9,0.9,0.9}{
            \makecell[{{p{\linewidth}}}]{
                \texttt{\tiny{[GM$\rangle$P2]}}
                \texttt{Edukalt teisaldatud 'A' asukohta (1,1). Teie uuendatud ruudustik näeb välja selline:} \\
\texttt{```} \\
\texttt{ 1234567} \\
\texttt{╔═╤═════╗} \\
\texttt{║A│S◌◌◌◌║ 1} \\
\texttt{╟─┤◌◌◌◌◌║ 2} \\
\texttt{║◌└─┐◌◌◌║ 3} \\
\texttt{║◌◌◌├─┐◌║ 4} \\
\texttt{║◌T◌│◌│◌║ 5} \\
\texttt{╟─┬─┴─┼─╢ 6} \\
\texttt{║◌│DMP│O║ 7} \\
\texttt{╚═╧═══╧═╝} \\
\texttt{```} \\
\texttt{Sa mängid praegu vooru 14 maksimaalsest 28 voorust.} \\
\texttt{Teil on kollektiivselt kogunenud 12 karistust 16 karistusest.} \\
\texttt{Teine mängija liigutas objekti oma ruudustikul.} \\
\\ 
\\ 
\texttt{Mis on sinu järgmine käsk?} \\
            }
        }
    }
    & \\ \\

    \theutterance \stepcounter{utterance}  
    & & & \multicolumn{4}{p{0.6\linewidth}}{
        \cellcolor[rgb]{0.9,0.9,0.9}{
            \makecell[{{p{\linewidth}}}]{
                \texttt{\tiny{[GM$\langle$P2]}}
                \texttt{LIIGUTA: D, (3, 1)} \\
            }
        }
    }
     \\ \\

    \theutterance \stepcounter{utterance}  
    & & & \multicolumn{2}{p{0.3\linewidth}}{
        \cellcolor[rgb]{0.9,0.9,0.9}{
            \makecell[{{p{\linewidth}}}]{
                \texttt{\tiny{[GM$|$GM]}}
                \texttt{LIIGUTA: D, (3, 1)} \\
            }
        }
    }
    & & \\ \\

    \theutterance \stepcounter{utterance}  
    & & & \multicolumn{2}{p{0.3\linewidth}}{
        \cellcolor[rgb]{0.9,0.9,0.9}{
            \makecell[{{p{\linewidth}}}]{
                \texttt{\tiny{[GM$|$GM]}}
                \texttt{Karistus: (3,1) ei ole tühi, vaid sisaldab objekti 'S'.} \\
\\ 
\texttt{Teil on kollektiivselt kogunenud 13 karistust 16 karistusest.} \\
\texttt{Palun proovi uuesti!} \\
            }
        }
    }
    & & \\ \\

    \theutterance \stepcounter{utterance}  
    & & \multicolumn{4}{p{0.6\linewidth}}{
        \cellcolor[rgb]{0.9,0.9,0.9}{
            \makecell[{{p{\linewidth}}}]{
                \texttt{\tiny{[GM$\rangle$P2]}}
                \texttt{Karistus: (3,1) ei ole tühi, vaid sisaldab objekti 'S'.} \\
\\ 
\texttt{Teil on kollektiivselt kogunenud 13 karistust 16 karistusest.} \\
\texttt{Palun proovi uuesti!} \\
            }
        }
    }
    & \\ \\

    \theutterance \stepcounter{utterance}  
    & & & \multicolumn{4}{p{0.6\linewidth}}{
        \cellcolor[rgb]{0.9,0.9,0.9}{
            \makecell[{{p{\linewidth}}}]{
                \texttt{\tiny{[GM$\langle$P2]}}
                \texttt{LIIGUTA: D, (3, 3)} \\
            }
        }
    }
     \\ \\

    \theutterance \stepcounter{utterance}  
    & & & \multicolumn{2}{p{0.3\linewidth}}{
        \cellcolor[rgb]{0.9,0.9,0.9}{
            \makecell[{{p{\linewidth}}}]{
                \texttt{\tiny{[GM$|$GM]}}
                \texttt{LIIGUTA: D, (3, 3)} \\
            }
        }
    }
    & & \\ \\

    \theutterance \stepcounter{utterance}  
    & & & \multicolumn{2}{p{0.3\linewidth}}{
        \cellcolor[rgb]{0.9,0.9,0.9}{
            \makecell[{{p{\linewidth}}}]{
                \texttt{\tiny{[GM$|$GM]}}
                \texttt{Karistus: (3,3) ei ole tühi, vaid sisaldab objekti '─'.} \\
\\ 
\texttt{Teil on kollektiivselt kogunenud 14 karistust 16 karistusest.} \\
\texttt{Palun proovi uuesti!} \\
            }
        }
    }
    & & \\ \\

    \theutterance \stepcounter{utterance}  
    & & \multicolumn{4}{p{0.6\linewidth}}{
        \cellcolor[rgb]{0.9,0.9,0.9}{
            \makecell[{{p{\linewidth}}}]{
                \texttt{\tiny{[GM$\rangle$P2]}}
                \texttt{Karistus: (3,3) ei ole tühi, vaid sisaldab objekti '─'.} \\
\\ 
\texttt{Teil on kollektiivselt kogunenud 14 karistust 16 karistusest.} \\
\texttt{Palun proovi uuesti!} \\
            }
        }
    }
    & \\ \\

    \theutterance \stepcounter{utterance}  
    & & & \multicolumn{4}{p{0.6\linewidth}}{
        \cellcolor[rgb]{0.9,0.9,0.9}{
            \makecell[{{p{\linewidth}}}]{
                \texttt{\tiny{[GM$\langle$P2]}}
                \texttt{LIIGUTA: D, (3, 4)} \\
            }
        }
    }
     \\ \\

    \theutterance \stepcounter{utterance}  
    & & & \multicolumn{2}{p{0.3\linewidth}}{
        \cellcolor[rgb]{0.9,0.9,0.9}{
            \makecell[{{p{\linewidth}}}]{
                \texttt{\tiny{[GM$|$GM]}}
                \texttt{LIIGUTA: D, (3, 4)} \\
            }
        }
    }
    & & \\ \\

    \theutterance \stepcounter{utterance}  
    & & & \multicolumn{2}{p{0.3\linewidth}}{
        \cellcolor[rgb]{0.9,0.9,0.9}{
            \makecell[{{p{\linewidth}}}]{
                \texttt{\tiny{[GM$|$GM]}}
                \texttt{D moved to (3, 4)} \\
            }
        }
    }
    & & \\ \\

    \theutterance \stepcounter{utterance}  
    & & \multicolumn{4}{p{0.6\linewidth}}{
        \cellcolor[rgb]{0.9,0.9,0.9}{
            \makecell[{{p{\linewidth}}}]{
                \texttt{\tiny{[P1$\langle$GM]}}
                \texttt{Edukalt teisaldatud 'D' asukohta (3,1). Teie uuendatud ruudustik näeb välja selline:} \\
\texttt{```} \\
\texttt{ 1234567} \\
\texttt{╔═╤═════╗} \\
\texttt{║A│D◌◌◌◌║ 1} \\
\texttt{╟─┤◌T◌◌◌║ 2} \\
\texttt{║◌└─┐◌◌◌║ 3} \\
\texttt{║◌◌◌├─┐◌║ 4} \\
\texttt{║◌◌S│◌│◌║ 5} \\
\texttt{╟─┬─┴─┼─╢ 6} \\
\texttt{║◌│O◌P│M║ 7} \\
\texttt{╚═╧═══╧═╝} \\
\texttt{```} \\
\texttt{Sa mängid praegu vooru 14 maksimaalsest 28 voorust.} \\
\texttt{Teil on kollektiivselt kogunenud 14 karistust 16 karistusest.} \\
\texttt{Teine mängija liigutas objekti oma ruudustikul.} \\
\\ 
\\ 
\texttt{Mis on sinu järgmine käsk?} \\
            }
        }
    }
    & \\ \\

    \theutterance \stepcounter{utterance}  
    & \multicolumn{4}{p{0.6\linewidth}}{
        \cellcolor[rgb]{0.9,0.9,0.9}{
            \makecell[{{p{\linewidth}}}]{
                \texttt{\tiny{[P1$\rangle$GM]}}
                \texttt{LIIGUTA: M, (4, 1)} \\
            }
        }
    }
    & & \\ \\

    \theutterance \stepcounter{utterance}  
    & & & \multicolumn{2}{p{0.3\linewidth}}{
        \cellcolor[rgb]{0.9,0.9,0.9}{
            \makecell[{{p{\linewidth}}}]{
                \texttt{\tiny{[GM$|$GM]}}
                \texttt{LIIGUTA: M, (4, 1)} \\
            }
        }
    }
    & & \\ \\

    \theutterance \stepcounter{utterance}  
    & & & \multicolumn{2}{p{0.3\linewidth}}{
        \cellcolor[rgb]{0.9,0.9,0.9}{
            \makecell[{{p{\linewidth}}}]{
                \texttt{\tiny{[GM$|$GM]}}
                \texttt{M moved to (4, 1)} \\
            }
        }
    }
    & & \\ \\

    \theutterance \stepcounter{utterance}  
    & & \multicolumn{4}{p{0.6\linewidth}}{
        \cellcolor[rgb]{0.9,0.9,0.9}{
            \makecell[{{p{\linewidth}}}]{
                \texttt{\tiny{[GM$\rangle$P2]}}
                \texttt{Edukalt teisaldatud 'D' asukohta (3,4). Teie uuendatud ruudustik näeb välja selline:} \\
\texttt{```} \\
\texttt{ 1234567} \\
\texttt{╔═╤═════╗} \\
\texttt{║A│S◌◌◌◌║ 1} \\
\texttt{╟─┤◌◌◌◌◌║ 2} \\
\texttt{║◌└─┐◌◌◌║ 3} \\
\texttt{║◌◌D├─┐◌║ 4} \\
\texttt{║◌T◌│◌│◌║ 5} \\
\texttt{╟─┬─┴─┼─╢ 6} \\
\texttt{║◌│◌MP│O║ 7} \\
\texttt{╚═╧═══╧═╝} \\
\texttt{```} \\
\texttt{Sa mängid praegu vooru 15 maksimaalsest 28 voorust.} \\
\texttt{Teil on kollektiivselt kogunenud 14 karistust 16 karistusest.} \\
\texttt{Teine mängija liigutas objekti oma ruudustikul.} \\
\\ 
\\ 
\texttt{Mis on sinu järgmine käsk?} \\
            }
        }
    }
    & \\ \\

    \theutterance \stepcounter{utterance}  
    & & & \multicolumn{4}{p{0.6\linewidth}}{
        \cellcolor[rgb]{0.9,0.9,0.9}{
            \makecell[{{p{\linewidth}}}]{
                \texttt{\tiny{[GM$\langle$P2]}}
                \texttt{LIIGUTA: M, (4, 1)} \\
            }
        }
    }
     \\ \\

    \theutterance \stepcounter{utterance}  
    & & & \multicolumn{2}{p{0.3\linewidth}}{
        \cellcolor[rgb]{0.9,0.9,0.9}{
            \makecell[{{p{\linewidth}}}]{
                \texttt{\tiny{[GM$|$GM]}}
                \texttt{LIIGUTA: M, (4, 1)} \\
            }
        }
    }
    & & \\ \\

    \theutterance \stepcounter{utterance}  
    & & & \multicolumn{2}{p{0.3\linewidth}}{
        \cellcolor[rgb]{0.9,0.9,0.9}{
            \makecell[{{p{\linewidth}}}]{
                \texttt{\tiny{[GM$|$GM]}}
                \texttt{M moved to (4, 1)} \\
            }
        }
    }
    & & \\ \\

    \theutterance \stepcounter{utterance}  
    & & \multicolumn{4}{p{0.6\linewidth}}{
        \cellcolor[rgb]{0.9,0.9,0.9}{
            \makecell[{{p{\linewidth}}}]{
                \texttt{\tiny{[P1$\langle$GM]}}
                \texttt{Edukalt teisaldatud 'M' asukohta (4,1). Teie uuendatud ruudustik näeb välja selline:} \\
\texttt{```} \\
\texttt{ 1234567} \\
\texttt{╔═╤═════╗} \\
\texttt{║A│DM◌◌◌║ 1} \\
\texttt{╟─┤◌T◌◌◌║ 2} \\
\texttt{║◌└─┐◌◌◌║ 3} \\
\texttt{║◌◌◌├─┐◌║ 4} \\
\texttt{║◌◌S│◌│◌║ 5} \\
\texttt{╟─┬─┴─┼─╢ 6} \\
\texttt{║◌│O◌P│◌║ 7} \\
\texttt{╚═╧═══╧═╝} \\
\texttt{```} \\
\texttt{Sa mängid praegu vooru 15 maksimaalsest 28 voorust.} \\
\texttt{Teil on kollektiivselt kogunenud 14 karistust 16 karistusest.} \\
\texttt{Teine mängija liigutas objekti oma ruudustikul.} \\
\\ 
\\ 
\texttt{Mis on sinu järgmine käsk?} \\
            }
        }
    }
    & \\ \\

    \theutterance \stepcounter{utterance}  
    & \multicolumn{4}{p{0.6\linewidth}}{
        \cellcolor[rgb]{0.9,0.9,0.9}{
            \makecell[{{p{\linewidth}}}]{
                \texttt{\tiny{[P1$\rangle$GM]}}
                \texttt{LIIGUTA: O, (5, 1)} \\
            }
        }
    }
    & & \\ \\

    \theutterance \stepcounter{utterance}  
    & & & \multicolumn{2}{p{0.3\linewidth}}{
        \cellcolor[rgb]{0.9,0.9,0.9}{
            \makecell[{{p{\linewidth}}}]{
                \texttt{\tiny{[GM$|$GM]}}
                \texttt{LIIGUTA: O, (5, 1)} \\
            }
        }
    }
    & & \\ \\

    \theutterance \stepcounter{utterance}  
    & & & \multicolumn{2}{p{0.3\linewidth}}{
        \cellcolor[rgb]{0.9,0.9,0.9}{
            \makecell[{{p{\linewidth}}}]{
                \texttt{\tiny{[GM$|$GM]}}
                \texttt{O moved to (5, 1)} \\
            }
        }
    }
    & & \\ \\

    \theutterance \stepcounter{utterance}  
    & & \multicolumn{4}{p{0.6\linewidth}}{
        \cellcolor[rgb]{0.9,0.9,0.9}{
            \makecell[{{p{\linewidth}}}]{
                \texttt{\tiny{[GM$\rangle$P2]}}
                \texttt{Edukalt teisaldatud 'M' asukohta (4,1). Teie uuendatud ruudustik näeb välja selline:} \\
\texttt{```} \\
\texttt{ 1234567} \\
\texttt{╔═╤═════╗} \\
\texttt{║A│SM◌◌◌║ 1} \\
\texttt{╟─┤◌◌◌◌◌║ 2} \\
\texttt{║◌└─┐◌◌◌║ 3} \\
\texttt{║◌◌D├─┐◌║ 4} \\
\texttt{║◌T◌│◌│◌║ 5} \\
\texttt{╟─┬─┴─┼─╢ 6} \\
\texttt{║◌│◌◌P│O║ 7} \\
\texttt{╚═╧═══╧═╝} \\
\texttt{```} \\
\texttt{Sa mängid praegu vooru 16 maksimaalsest 28 voorust.} \\
\texttt{Teil on kollektiivselt kogunenud 14 karistust 16 karistusest.} \\
\texttt{Teine mängija liigutas objekti oma ruudustikul.} \\
\\ 
\\ 
\texttt{Mis on sinu järgmine käsk?} \\
            }
        }
    }
    & \\ \\

    \theutterance \stepcounter{utterance}  
    & & & \multicolumn{4}{p{0.6\linewidth}}{
        \cellcolor[rgb]{0.9,0.9,0.9}{
            \makecell[{{p{\linewidth}}}]{
                \texttt{\tiny{[GM$\langle$P2]}}
                \texttt{LIIGUTA: O, (5, 1)} \\
            }
        }
    }
     \\ \\

    \theutterance \stepcounter{utterance}  
    & & & \multicolumn{2}{p{0.3\linewidth}}{
        \cellcolor[rgb]{0.9,0.9,0.9}{
            \makecell[{{p{\linewidth}}}]{
                \texttt{\tiny{[GM$|$GM]}}
                \texttt{LIIGUTA: O, (5, 1)} \\
            }
        }
    }
    & & \\ \\

    \theutterance \stepcounter{utterance}  
    & & & \multicolumn{2}{p{0.3\linewidth}}{
        \cellcolor[rgb]{0.9,0.9,0.9}{
            \makecell[{{p{\linewidth}}}]{
                \texttt{\tiny{[GM$|$GM]}}
                \texttt{O moved to (5, 1)} \\
            }
        }
    }
    & & \\ \\

    \theutterance \stepcounter{utterance}  
    & & \multicolumn{4}{p{0.6\linewidth}}{
        \cellcolor[rgb]{0.9,0.9,0.9}{
            \makecell[{{p{\linewidth}}}]{
                \texttt{\tiny{[P1$\langle$GM]}}
                \texttt{Edukalt teisaldatud 'O' asukohta (5,1). Teie uuendatud ruudustik näeb välja selline:} \\
\texttt{```} \\
\texttt{ 1234567} \\
\texttt{╔═╤═════╗} \\
\texttt{║A│DMO◌◌║ 1} \\
\texttt{╟─┤◌T◌◌◌║ 2} \\
\texttt{║◌└─┐◌◌◌║ 3} \\
\texttt{║◌◌◌├─┐◌║ 4} \\
\texttt{║◌◌S│◌│◌║ 5} \\
\texttt{╟─┬─┴─┼─╢ 6} \\
\texttt{║◌│◌◌P│◌║ 7} \\
\texttt{╚═╧═══╧═╝} \\
\texttt{```} \\
\texttt{Sa mängid praegu vooru 16 maksimaalsest 28 voorust.} \\
\texttt{Teil on kollektiivselt kogunenud 14 karistust 16 karistusest.} \\
\texttt{Teine mängija liigutas objekti oma ruudustikul.} \\
\\ 
\\ 
\texttt{Mis on sinu järgmine käsk?} \\
            }
        }
    }
    & \\ \\

    \theutterance \stepcounter{utterance}  
    & \multicolumn{4}{p{0.6\linewidth}}{
        \cellcolor[rgb]{0.9,0.9,0.9}{
            \makecell[{{p{\linewidth}}}]{
                \texttt{\tiny{[P1$\rangle$GM]}}
                \texttt{LIIGUTA: P, (6, 1)} \\
            }
        }
    }
    & & \\ \\

    \theutterance \stepcounter{utterance}  
    & & & \multicolumn{2}{p{0.3\linewidth}}{
        \cellcolor[rgb]{0.9,0.9,0.9}{
            \makecell[{{p{\linewidth}}}]{
                \texttt{\tiny{[GM$|$GM]}}
                \texttt{LIIGUTA: P, (6, 1)} \\
            }
        }
    }
    & & \\ \\

    \theutterance \stepcounter{utterance}  
    & & & \multicolumn{2}{p{0.3\linewidth}}{
        \cellcolor[rgb]{0.9,0.9,0.9}{
            \makecell[{{p{\linewidth}}}]{
                \texttt{\tiny{[GM$|$GM]}}
                \texttt{P moved to (6, 1)} \\
            }
        }
    }
    & & \\ \\

    \theutterance \stepcounter{utterance}  
    & & \multicolumn{4}{p{0.6\linewidth}}{
        \cellcolor[rgb]{0.9,0.9,0.9}{
            \makecell[{{p{\linewidth}}}]{
                \texttt{\tiny{[GM$\rangle$P2]}}
                \texttt{Edukalt teisaldatud 'O' asukohta (5,1). Teie uuendatud ruudustik näeb välja selline:} \\
\texttt{```} \\
\texttt{ 1234567} \\
\texttt{╔═╤═════╗} \\
\texttt{║A│SMO◌◌║ 1} \\
\texttt{╟─┤◌◌◌◌◌║ 2} \\
\texttt{║◌└─┐◌◌◌║ 3} \\
\texttt{║◌◌D├─┐◌║ 4} \\
\texttt{║◌T◌│◌│◌║ 5} \\
\texttt{╟─┬─┴─┼─╢ 6} \\
\texttt{║◌│◌◌P│◌║ 7} \\
\texttt{╚═╧═══╧═╝} \\
\texttt{```} \\
\texttt{Sa mängid praegu vooru 17 maksimaalsest 28 voorust.} \\
\texttt{Teil on kollektiivselt kogunenud 14 karistust 16 karistusest.} \\
\texttt{Teine mängija liigutas objekti oma ruudustikul.} \\
\\ 
\\ 
\texttt{Mis on sinu järgmine käsk?} \\
            }
        }
    }
    & \\ \\

    \theutterance \stepcounter{utterance}  
    & & & \multicolumn{4}{p{0.6\linewidth}}{
        \cellcolor[rgb]{0.9,0.9,0.9}{
            \makecell[{{p{\linewidth}}}]{
                \texttt{\tiny{[GM$\langle$P2]}}
                \texttt{LIIGUTA: P, (6, 1)} \\
            }
        }
    }
     \\ \\

    \theutterance \stepcounter{utterance}  
    & & & \multicolumn{2}{p{0.3\linewidth}}{
        \cellcolor[rgb]{0.9,0.9,0.9}{
            \makecell[{{p{\linewidth}}}]{
                \texttt{\tiny{[GM$|$GM]}}
                \texttt{LIIGUTA: P, (6, 1)} \\
            }
        }
    }
    & & \\ \\

    \theutterance \stepcounter{utterance}  
    & & & \multicolumn{2}{p{0.3\linewidth}}{
        \cellcolor[rgb]{0.9,0.9,0.9}{
            \makecell[{{p{\linewidth}}}]{
                \texttt{\tiny{[GM$|$GM]}}
                \texttt{P moved to (6, 1)} \\
            }
        }
    }
    & & \\ \\

    \theutterance \stepcounter{utterance}  
    & & \multicolumn{4}{p{0.6\linewidth}}{
        \cellcolor[rgb]{0.9,0.9,0.9}{
            \makecell[{{p{\linewidth}}}]{
                \texttt{\tiny{[P1$\langle$GM]}}
                \texttt{Edukalt teisaldatud 'P' asukohta (6,1). Teie uuendatud ruudustik näeb välja selline:} \\
\texttt{```} \\
\texttt{ 1234567} \\
\texttt{╔═╤═════╗} \\
\texttt{║A│DMOP◌║ 1} \\
\texttt{╟─┤◌T◌◌◌║ 2} \\
\texttt{║◌└─┐◌◌◌║ 3} \\
\texttt{║◌◌◌├─┐◌║ 4} \\
\texttt{║◌◌S│◌│◌║ 5} \\
\texttt{╟─┬─┴─┼─╢ 6} \\
\texttt{║◌│◌◌◌│◌║ 7} \\
\texttt{╚═╧═══╧═╝} \\
\texttt{```} \\
\texttt{Sa mängid praegu vooru 17 maksimaalsest 28 voorust.} \\
\texttt{Teil on kollektiivselt kogunenud 14 karistust 16 karistusest.} \\
\texttt{Teine mängija liigutas objekti oma ruudustikul.} \\
\\ 
\\ 
\texttt{Mis on sinu järgmine käsk?} \\
            }
        }
    }
    & \\ \\

    \theutterance \stepcounter{utterance}  
    & \multicolumn{4}{p{0.6\linewidth}}{
        \cellcolor[rgb]{0.9,0.9,0.9}{
            \makecell[{{p{\linewidth}}}]{
                \texttt{\tiny{[P1$\rangle$GM]}}
                \texttt{LIIGUTA: S, (7, 1)} \\
            }
        }
    }
    & & \\ \\

    \theutterance \stepcounter{utterance}  
    & & & \multicolumn{2}{p{0.3\linewidth}}{
        \cellcolor[rgb]{0.9,0.9,0.9}{
            \makecell[{{p{\linewidth}}}]{
                \texttt{\tiny{[GM$|$GM]}}
                \texttt{LIIGUTA: S, (7, 1)} \\
            }
        }
    }
    & & \\ \\

    \theutterance \stepcounter{utterance}  
    & & & \multicolumn{2}{p{0.3\linewidth}}{
        \cellcolor[rgb]{0.9,0.9,0.9}{
            \makecell[{{p{\linewidth}}}]{
                \texttt{\tiny{[GM$|$GM]}}
                \texttt{S moved to (7, 1)} \\
            }
        }
    }
    & & \\ \\

    \theutterance \stepcounter{utterance}  
    & & \multicolumn{4}{p{0.6\linewidth}}{
        \cellcolor[rgb]{0.9,0.9,0.9}{
            \makecell[{{p{\linewidth}}}]{
                \texttt{\tiny{[GM$\rangle$P2]}}
                \texttt{Edukalt teisaldatud 'P' asukohta (6,1). Teie uuendatud ruudustik näeb välja selline:} \\
\texttt{```} \\
\texttt{ 1234567} \\
\texttt{╔═╤═════╗} \\
\texttt{║A│SMOP◌║ 1} \\
\texttt{╟─┤◌◌◌◌◌║ 2} \\
\texttt{║◌└─┐◌◌◌║ 3} \\
\texttt{║◌◌D├─┐◌║ 4} \\
\texttt{║◌T◌│◌│◌║ 5} \\
\texttt{╟─┬─┴─┼─╢ 6} \\
\texttt{║◌│◌◌◌│◌║ 7} \\
\texttt{╚═╧═══╧═╝} \\
\texttt{```} \\
\texttt{Sa mängid praegu vooru 18 maksimaalsest 28 voorust.} \\
\texttt{Teil on kollektiivselt kogunenud 14 karistust 16 karistusest.} \\
\texttt{Teine mängija liigutas objekti oma ruudustikul.} \\
\\ 
\\ 
\texttt{Mis on sinu järgmine käsk?} \\
            }
        }
    }
    & \\ \\

    \theutterance \stepcounter{utterance}  
    & & & \multicolumn{4}{p{0.6\linewidth}}{
        \cellcolor[rgb]{0.9,0.9,0.9}{
            \makecell[{{p{\linewidth}}}]{
                \texttt{\tiny{[GM$\langle$P2]}}
                \texttt{LIIGUTA: S, (7, 1)} \\
            }
        }
    }
     \\ \\

    \theutterance \stepcounter{utterance}  
    & & & \multicolumn{2}{p{0.3\linewidth}}{
        \cellcolor[rgb]{0.9,0.9,0.9}{
            \makecell[{{p{\linewidth}}}]{
                \texttt{\tiny{[GM$|$GM]}}
                \texttt{LIIGUTA: S, (7, 1)} \\
            }
        }
    }
    & & \\ \\

    \theutterance \stepcounter{utterance}  
    & & & \multicolumn{2}{p{0.3\linewidth}}{
        \cellcolor[rgb]{0.9,0.9,0.9}{
            \makecell[{{p{\linewidth}}}]{
                \texttt{\tiny{[GM$|$GM]}}
                \texttt{S moved to (7, 1)} \\
            }
        }
    }
    & & \\ \\

    \theutterance \stepcounter{utterance}  
    & & \multicolumn{4}{p{0.6\linewidth}}{
        \cellcolor[rgb]{0.9,0.9,0.9}{
            \makecell[{{p{\linewidth}}}]{
                \texttt{\tiny{[P1$\langle$GM]}}
                \texttt{Edukalt teisaldatud 'S' asukohta (7,1). Teie uuendatud ruudustik näeb välja selline:} \\
\texttt{```} \\
\texttt{ 1234567} \\
\texttt{╔═╤═════╗} \\
\texttt{║A│DMOPS║ 1} \\
\texttt{╟─┤◌T◌◌◌║ 2} \\
\texttt{║◌└─┐◌◌◌║ 3} \\
\texttt{║◌◌◌├─┐◌║ 4} \\
\texttt{║◌◌◌│◌│◌║ 5} \\
\texttt{╟─┬─┴─┼─╢ 6} \\
\texttt{║◌│◌◌◌│◌║ 7} \\
\texttt{╚═╧═══╧═╝} \\
\texttt{```} \\
\texttt{Sa mängid praegu vooru 18 maksimaalsest 28 voorust.} \\
\texttt{Teil on kollektiivselt kogunenud 14 karistust 16 karistusest.} \\
\texttt{Teine mängija liigutas objekti oma ruudustikul.} \\
\\ 
\\ 
\texttt{Mis on sinu järgmine käsk?} \\
            }
        }
    }
    & \\ \\

    \theutterance \stepcounter{utterance}  
    & \multicolumn{4}{p{0.6\linewidth}}{
        \cellcolor[rgb]{0.9,0.9,0.9}{
            \makecell[{{p{\linewidth}}}]{
                \texttt{\tiny{[P1$\rangle$GM]}}
                \texttt{LIIGUTA: T, (2, 4)} \\
            }
        }
    }
    & & \\ \\

    \theutterance \stepcounter{utterance}  
    & & & \multicolumn{2}{p{0.3\linewidth}}{
        \cellcolor[rgb]{0.9,0.9,0.9}{
            \makecell[{{p{\linewidth}}}]{
                \texttt{\tiny{[GM$|$GM]}}
                \texttt{LIIGUTA: T, (2, 4)} \\
            }
        }
    }
    & & \\ \\

    \theutterance \stepcounter{utterance}  
    & & & \multicolumn{2}{p{0.3\linewidth}}{
        \cellcolor[rgb]{0.9,0.9,0.9}{
            \makecell[{{p{\linewidth}}}]{
                \texttt{\tiny{[GM$|$GM]}}
                \texttt{T moved to (2, 4)} \\
            }
        }
    }
    & & \\ \\

    \theutterance \stepcounter{utterance}  
    & & \multicolumn{4}{p{0.6\linewidth}}{
        \cellcolor[rgb]{0.9,0.9,0.9}{
            \makecell[{{p{\linewidth}}}]{
                \texttt{\tiny{[GM$\rangle$P2]}}
                \texttt{Edukalt teisaldatud 'S' asukohta (7,1). Teie uuendatud ruudustik näeb välja selline:} \\
\texttt{```} \\
\texttt{ 1234567} \\
\texttt{╔═╤═════╗} \\
\texttt{║A│◌MOPS║ 1} \\
\texttt{╟─┤◌◌◌◌◌║ 2} \\
\texttt{║◌└─┐◌◌◌║ 3} \\
\texttt{║◌◌D├─┐◌║ 4} \\
\texttt{║◌T◌│◌│◌║ 5} \\
\texttt{╟─┬─┴─┼─╢ 6} \\
\texttt{║◌│◌◌◌│◌║ 7} \\
\texttt{╚═╧═══╧═╝} \\
\texttt{```} \\
\texttt{Sa mängid praegu vooru 19 maksimaalsest 28 voorust.} \\
\texttt{Teil on kollektiivselt kogunenud 14 karistust 16 karistusest.} \\
\texttt{Teine mängija liigutas objekti oma ruudustikul.} \\
\\ 
\\ 
\texttt{Mis on sinu järgmine käsk?} \\
            }
        }
    }
    & \\ \\

    \theutterance \stepcounter{utterance}  
    & & & \multicolumn{4}{p{0.6\linewidth}}{
        \cellcolor[rgb]{0.9,0.9,0.9}{
            \makecell[{{p{\linewidth}}}]{
                \texttt{\tiny{[GM$\langle$P2]}}
                \texttt{LIIGUTA: T, (2, 4)} \\
            }
        }
    }
     \\ \\

    \theutterance \stepcounter{utterance}  
    & & & \multicolumn{2}{p{0.3\linewidth}}{
        \cellcolor[rgb]{0.9,0.9,0.9}{
            \makecell[{{p{\linewidth}}}]{
                \texttt{\tiny{[GM$|$GM]}}
                \texttt{LIIGUTA: T, (2, 4)} \\
            }
        }
    }
    & & \\ \\

    \theutterance \stepcounter{utterance}  
    & & & \multicolumn{2}{p{0.3\linewidth}}{
        \cellcolor[rgb]{0.9,0.9,0.9}{
            \makecell[{{p{\linewidth}}}]{
                \texttt{\tiny{[GM$|$GM]}}
                \texttt{T moved to (2, 4)} \\
            }
        }
    }
    & & \\ \\

    \theutterance \stepcounter{utterance}  
    & & \multicolumn{4}{p{0.6\linewidth}}{
        \cellcolor[rgb]{0.9,0.9,0.9}{
            \makecell[{{p{\linewidth}}}]{
                \texttt{\tiny{[P1$\langle$GM]}}
                \texttt{Edukalt teisaldatud 'T' asukohta (2,4). Teie uuendatud ruudustik näeb välja selline:} \\
\texttt{```} \\
\texttt{ 1234567} \\
\texttt{╔═╤═════╗} \\
\texttt{║A│DMOPS║ 1} \\
\texttt{╟─┤◌◌◌◌◌║ 2} \\
\texttt{║◌└─┐◌◌◌║ 3} \\
\texttt{║◌T◌├─┐◌║ 4} \\
\texttt{║◌◌◌│◌│◌║ 5} \\
\texttt{╟─┬─┴─┼─╢ 6} \\
\texttt{║◌│◌◌◌│◌║ 7} \\
\texttt{╚═╧═══╧═╝} \\
\texttt{```} \\
\texttt{Sa mängid praegu vooru 19 maksimaalsest 28 voorust.} \\
\texttt{Teil on kollektiivselt kogunenud 14 karistust 16 karistusest.} \\
\texttt{Teine mängija liigutas objekti oma ruudustikul.} \\
\\ 
\\ 
\texttt{Mis on sinu järgmine käsk?} \\
            }
        }
    }
    & \\ \\

    \theutterance \stepcounter{utterance}  
    & \multicolumn{4}{p{0.6\linewidth}}{
        \cellcolor[rgb]{0.9,0.9,0.9}{
            \makecell[{{p{\linewidth}}}]{
                \texttt{\tiny{[P1$\rangle$GM]}}
                \texttt{ÜTLE: Kõik minu 7 objekti on sihtkoordinaatides: A(1,1), D(3,1), M(4,1), O(5,1), P(6,1), S(7,1), T(2,4). Kas teie laual on kõik samades asendites? Kui jah, siis valmis!} \\
            }
        }
    }
    & & \\ \\

    \theutterance \stepcounter{utterance}  
    & & & \multicolumn{2}{p{0.3\linewidth}}{
        \cellcolor[rgb]{0.9,0.9,0.9}{
            \makecell[{{p{\linewidth}}}]{
                \texttt{\tiny{[GM$|$GM]}}
                \texttt{ÜTLE: Kõik minu 7 objekti on sihtkoordinaatides: A(1,1), D(3,1), M(4,1), O(5,1), P(6,1), S(7,1), T(2,4). Kas teie laual on kõik samades asendites? Kui jah, siis valmis!} \\
            }
        }
    }
    & & \\ \\

    \theutterance \stepcounter{utterance}  
    & & \multicolumn{4}{p{0.6\linewidth}}{
        \cellcolor[rgb]{0.9,0.9,0.9}{
            \makecell[{{p{\linewidth}}}]{
                \texttt{\tiny{[GM$\rangle$P2]}}
                \texttt{Edukalt teisaldatud 'T' asukohta (2,4). Teie uuendatud ruudustik näeb välja selline:} \\
\texttt{```} \\
\texttt{ 1234567} \\
\texttt{╔═╤═════╗} \\
\texttt{║A│◌MOPS║ 1} \\
\texttt{╟─┤◌◌◌◌◌║ 2} \\
\texttt{║◌└─┐◌◌◌║ 3} \\
\texttt{║◌TD├─┐◌║ 4} \\
\texttt{║◌◌◌│◌│◌║ 5} \\
\texttt{╟─┬─┴─┼─╢ 6} \\
\texttt{║◌│◌◌◌│◌║ 7} \\
\texttt{╚═╧═══╧═╝} \\
\texttt{```} \\
\texttt{Sa mängid praegu vooru 20 maksimaalsest 28 voorust.} \\
\texttt{Teil on kollektiivselt kogunenud 14 karistust 16 karistusest.} \\
\texttt{Teine mängija saatis selle sõnumi:} \\
\\ 
\\ 
\texttt{"Kõik minu 7 objekti on sihtkoordinaatides: A(1,1), D(3,1), M(4,1), O(5,1), P(6,1), S(7,1), T(2,4). Kas teie laual on kõik samades asendites? Kui jah, siis valmis!"} \\
\\ 
\\ 
\texttt{Mis on sinu järgmine käsk?} \\
            }
        }
    }
    & \\ \\

    \theutterance \stepcounter{utterance}  
    & & & \multicolumn{4}{p{0.6\linewidth}}{
        \cellcolor[rgb]{0.9,0.9,0.9}{
            \makecell[{{p{\linewidth}}}]{
                \texttt{\tiny{[GM$\langle$P2]}}
                \texttt{ÜTLE: valmis!} \\
            }
        }
    }
     \\ \\

    \theutterance \stepcounter{utterance}  
    & & & \multicolumn{2}{p{0.3\linewidth}}{
        \cellcolor[rgb]{0.9,0.9,0.9}{
            \makecell[{{p{\linewidth}}}]{
                \texttt{\tiny{[GM$|$GM]}}
                \texttt{ÜTLE: valmis!} \\
            }
        }
    }
    & & \\ \\

    \theutterance \stepcounter{utterance}  
    & & & \multicolumn{2}{p{0.3\linewidth}}{
        \cellcolor[rgb]{0.9,0.9,0.9}{
            \makecell[{{p{\linewidth}}}]{
                \texttt{\tiny{[GM$|$GM]}}
                \texttt{true} \\
            }
        }
    }
    & & \\ \\

    \theutterance \stepcounter{utterance}  
    & & & \multicolumn{2}{p{0.3\linewidth}}{
        \cellcolor[rgb]{0.9,0.9,0.9}{
            \makecell[{{p{\linewidth}}}]{
                \texttt{\tiny{[GM$|$GM]}}
                \texttt{Initial states:} \\
\texttt{```} \\
\texttt{Player 1:} \\
\texttt{ 1234567} \\
\texttt{╔═╤═════╗} \\
\texttt{║◌│◌A◌◌◌║ 1} \\
\texttt{╟─┤◌T◌◌◌║ 2} \\
\texttt{║O└─┐◌◌◌║ 3} \\
\texttt{║◌◌◌├─┐◌║ 4} \\
\texttt{║◌◌S│◌│◌║ 5} \\
\texttt{╟─┬─┴─┼─╢ 6} \\
\texttt{║P│◌DM│◌║ 7} \\
\texttt{╚═╧═══╧═╝} \\
\\ 
\\ 
\texttt{Player 2:} \\
\texttt{ 1234567} \\
\texttt{╔═╤═════╗} \\
\texttt{║P│S◌D◌◌║ 1} \\
\texttt{╟─┤◌◌◌◌◌║ 2} \\
\texttt{║◌└─┐◌A◌║ 3} \\
\texttt{║◌◌M├─┐◌║ 4} \\
\texttt{║◌T◌│◌│O║ 5} \\
\texttt{╟─┬─┴─┼─╢ 6} \\
\texttt{║◌│◌◌◌│◌║ 7} \\
\texttt{╚═╧═══╧═╝} \\
\\ 
\texttt{```} \\
            }
        }
    }
    & & \\ \\

    \theutterance \stepcounter{utterance}  
    & & & \multicolumn{2}{p{0.3\linewidth}}{
        \cellcolor[rgb]{0.9,0.9,0.9}{
            \makecell[{{p{\linewidth}}}]{
                \texttt{\tiny{[GM$|$GM]}}
                \texttt{End states:} \\
\texttt{```} \\
\texttt{Player 1:} \\
\texttt{ 1234567} \\
\texttt{╔═╤═════╗} \\
\texttt{║A│DMOPS║ 1} \\
\texttt{╟─┤◌◌◌◌◌║ 2} \\
\texttt{║◌└─┐◌◌◌║ 3} \\
\texttt{║◌T◌├─┐◌║ 4} \\
\texttt{║◌◌◌│◌│◌║ 5} \\
\texttt{╟─┬─┴─┼─╢ 6} \\
\texttt{║◌│◌◌◌│◌║ 7} \\
\texttt{╚═╧═══╧═╝} \\
\\ 
\\ 
\texttt{Player 2:} \\
\texttt{ 1234567} \\
\texttt{╔═╤═════╗} \\
\texttt{║A│◌MOPS║ 1} \\
\texttt{╟─┤◌◌◌◌◌║ 2} \\
\texttt{║◌└─┐◌◌◌║ 3} \\
\texttt{║◌TD├─┐◌║ 4} \\
\texttt{║◌◌◌│◌│◌║ 5} \\
\texttt{╟─┬─┴─┼─╢ 6} \\
\texttt{║◌│◌◌◌│◌║ 7} \\
\texttt{╚═╧═══╧═╝} \\
\\ 
\texttt{```} \\
            }
        }
    }
    & & \\ \\

    \theutterance \stepcounter{utterance}  
    & & & \multicolumn{2}{p{0.3\linewidth}}{
        \cellcolor[rgb]{0.9,0.9,0.9}{
            \makecell[{{p{\linewidth}}}]{
                \texttt{\tiny{[GM$|$GM]}}
                \texttt{* success: True} \\
\texttt{* lose: False} \\
\texttt{* aborted: False} \\
\texttt{{-}{-}{-}{-}{-}{-}{-}} \\
\texttt{* Shifts: 14.00} \\
\texttt{* Max Shifts: 12.00} \\
\texttt{* Min Shifts: 6.00} \\
\texttt{* End Distance Sum: 3.00} \\
\texttt{* Init Distance Sum: 32.45} \\
\texttt{* Expected Distance Sum: 29.33} \\
\texttt{* Penalties: 14.00} \\
\texttt{* Max Penalties: 16.00} \\
\texttt{* Rounds: 20.00} \\
\texttt{* Max Rounds: 28.00} \\
\texttt{* Object Count: 7.00} \\
            }
        }
    }
    & & \\ \\

\end{supertabular}
}

\end{document}
